\documentclass[11pt]{article}

    \usepackage[breakable]{tcolorbox}
    \usepackage{parskip} % Stop auto-indenting (to mimic markdown behaviour)
    

    % Basic figure setup, for now with no caption control since it's done
    % automatically by Pandoc (which extracts ![](path) syntax from Markdown).
    \usepackage{graphicx}
    % Keep aspect ratio if custom image width or height is specified
    \setkeys{Gin}{keepaspectratio}
    % Maintain compatibility with old templates. Remove in nbconvert 6.0
    \let\Oldincludegraphics\includegraphics
    % Ensure that by default, figures have no caption (until we provide a
    % proper Figure object with a Caption API and a way to capture that
    % in the conversion process - todo).
    \usepackage{caption}
    \DeclareCaptionFormat{nocaption}{}
    \captionsetup{format=nocaption,aboveskip=0pt,belowskip=0pt}

    \usepackage{float}
    \floatplacement{figure}{H} % forces figures to be placed at the correct location
    \usepackage{xcolor} % Allow colors to be defined
    \usepackage{enumerate} % Needed for markdown enumerations to work
    \usepackage{geometry} % Used to adjust the document margins
    \usepackage{amsmath} % Equations
    \usepackage{amssymb} % Equations
    \usepackage{textcomp} % defines textquotesingle
    % Hack from http://tex.stackexchange.com/a/47451/13684:
    \AtBeginDocument{%
        \def\PYZsq{\textquotesingle}% Upright quotes in Pygmentized code
    }
    \usepackage{upquote} % Upright quotes for verbatim code
    \usepackage{eurosym} % defines \euro

    \usepackage{iftex}
    \ifPDFTeX
        \usepackage[T1]{fontenc}
        \IfFileExists{alphabeta.sty}{
              \usepackage{alphabeta}
          }{
              \usepackage[mathletters]{ucs}
              \usepackage[utf8x]{inputenc}
          }
    \else
        \usepackage{fontspec}
        \usepackage{unicode-math}
    \fi

    \usepackage{fancyvrb} % verbatim replacement that allows latex
    \usepackage{grffile} % extends the file name processing of package graphics
                         % to support a larger range
    \makeatletter % fix for old versions of grffile with XeLaTeX
    \@ifpackagelater{grffile}{2019/11/01}
    {
      % Do nothing on new versions
    }
    {
      \def\Gread@@xetex#1{%
        \IfFileExists{"\Gin@base".bb}%
        {\Gread@eps{\Gin@base.bb}}%
        {\Gread@@xetex@aux#1}%
      }
    }
    \makeatother
    \usepackage[Export]{adjustbox} % Used to constrain images to a maximum size
    \adjustboxset{max size={0.9\linewidth}{0.9\paperheight}}

    % The hyperref package gives us a pdf with properly built
    % internal navigation ('pdf bookmarks' for the table of contents,
    % internal cross-reference links, web links for URLs, etc.)
    \usepackage{hyperref}
    % The default LaTeX title has an obnoxious amount of whitespace. By default,
    % titling removes some of it. It also provides customization options.
    \usepackage{titling}
    \usepackage{longtable} % longtable support required by pandoc >1.10
    \usepackage{booktabs}  % table support for pandoc > 1.12.2
    \usepackage{array}     % table support for pandoc >= 2.11.3
    \usepackage{calc}      % table minipage width calculation for pandoc >= 2.11.1
    \usepackage[inline]{enumitem} % IRkernel/repr support (it uses the enumerate* environment)
    \usepackage[normalem]{ulem} % ulem is needed to support strikethroughs (\sout)
                                % normalem makes italics be italics, not underlines
    \usepackage{soul}      % strikethrough (\st) support for pandoc >= 3.0.0
    \usepackage{mathrsfs}
    

    
    % Colors for the hyperref package
    \definecolor{urlcolor}{rgb}{0,.145,.698}
    \definecolor{linkcolor}{rgb}{.71,0.21,0.01}
    \definecolor{citecolor}{rgb}{.12,.54,.11}

    % ANSI colors
    \definecolor{ansi-black}{HTML}{3E424D}
    \definecolor{ansi-black-intense}{HTML}{282C36}
    \definecolor{ansi-red}{HTML}{E75C58}
    \definecolor{ansi-red-intense}{HTML}{B22B31}
    \definecolor{ansi-green}{HTML}{00A250}
    \definecolor{ansi-green-intense}{HTML}{007427}
    \definecolor{ansi-yellow}{HTML}{DDB62B}
    \definecolor{ansi-yellow-intense}{HTML}{B27D12}
    \definecolor{ansi-blue}{HTML}{208FFB}
    \definecolor{ansi-blue-intense}{HTML}{0065CA}
    \definecolor{ansi-magenta}{HTML}{D160C4}
    \definecolor{ansi-magenta-intense}{HTML}{A03196}
    \definecolor{ansi-cyan}{HTML}{60C6C8}
    \definecolor{ansi-cyan-intense}{HTML}{258F8F}
    \definecolor{ansi-white}{HTML}{C5C1B4}
    \definecolor{ansi-white-intense}{HTML}{A1A6B2}
    \definecolor{ansi-default-inverse-fg}{HTML}{FFFFFF}
    \definecolor{ansi-default-inverse-bg}{HTML}{000000}

    % common color for the border for error outputs.
    \definecolor{outerrorbackground}{HTML}{FFDFDF}

    % commands and environments needed by pandoc snippets
    % extracted from the output of `pandoc -s`
    \providecommand{\tightlist}{%
      \setlength{\itemsep}{0pt}\setlength{\parskip}{0pt}}
    \DefineVerbatimEnvironment{Highlighting}{Verbatim}{commandchars=\\\{\}}
    % Add ',fontsize=\small' for more characters per line
    \newenvironment{Shaded}{}{}
    \newcommand{\KeywordTok}[1]{\textcolor[rgb]{0.00,0.44,0.13}{\textbf{{#1}}}}
    \newcommand{\DataTypeTok}[1]{\textcolor[rgb]{0.56,0.13,0.00}{{#1}}}
    \newcommand{\DecValTok}[1]{\textcolor[rgb]{0.25,0.63,0.44}{{#1}}}
    \newcommand{\BaseNTok}[1]{\textcolor[rgb]{0.25,0.63,0.44}{{#1}}}
    \newcommand{\FloatTok}[1]{\textcolor[rgb]{0.25,0.63,0.44}{{#1}}}
    \newcommand{\CharTok}[1]{\textcolor[rgb]{0.25,0.44,0.63}{{#1}}}
    \newcommand{\StringTok}[1]{\textcolor[rgb]{0.25,0.44,0.63}{{#1}}}
    \newcommand{\CommentTok}[1]{\textcolor[rgb]{0.38,0.63,0.69}{\textit{{#1}}}}
    \newcommand{\OtherTok}[1]{\textcolor[rgb]{0.00,0.44,0.13}{{#1}}}
    \newcommand{\AlertTok}[1]{\textcolor[rgb]{1.00,0.00,0.00}{\textbf{{#1}}}}
    \newcommand{\FunctionTok}[1]{\textcolor[rgb]{0.02,0.16,0.49}{{#1}}}
    \newcommand{\RegionMarkerTok}[1]{{#1}}
    \newcommand{\ErrorTok}[1]{\textcolor[rgb]{1.00,0.00,0.00}{\textbf{{#1}}}}
    \newcommand{\NormalTok}[1]{{#1}}

    % Additional commands for more recent versions of Pandoc
    \newcommand{\ConstantTok}[1]{\textcolor[rgb]{0.53,0.00,0.00}{{#1}}}
    \newcommand{\SpecialCharTok}[1]{\textcolor[rgb]{0.25,0.44,0.63}{{#1}}}
    \newcommand{\VerbatimStringTok}[1]{\textcolor[rgb]{0.25,0.44,0.63}{{#1}}}
    \newcommand{\SpecialStringTok}[1]{\textcolor[rgb]{0.73,0.40,0.53}{{#1}}}
    \newcommand{\ImportTok}[1]{{#1}}
    \newcommand{\DocumentationTok}[1]{\textcolor[rgb]{0.73,0.13,0.13}{\textit{{#1}}}}
    \newcommand{\AnnotationTok}[1]{\textcolor[rgb]{0.38,0.63,0.69}{\textbf{\textit{{#1}}}}}
    \newcommand{\CommentVarTok}[1]{\textcolor[rgb]{0.38,0.63,0.69}{\textbf{\textit{{#1}}}}}
    \newcommand{\VariableTok}[1]{\textcolor[rgb]{0.10,0.09,0.49}{{#1}}}
    \newcommand{\ControlFlowTok}[1]{\textcolor[rgb]{0.00,0.44,0.13}{\textbf{{#1}}}}
    \newcommand{\OperatorTok}[1]{\textcolor[rgb]{0.40,0.40,0.40}{{#1}}}
    \newcommand{\BuiltInTok}[1]{{#1}}
    \newcommand{\ExtensionTok}[1]{{#1}}
    \newcommand{\PreprocessorTok}[1]{\textcolor[rgb]{0.74,0.48,0.00}{{#1}}}
    \newcommand{\AttributeTok}[1]{\textcolor[rgb]{0.49,0.56,0.16}{{#1}}}
    \newcommand{\InformationTok}[1]{\textcolor[rgb]{0.38,0.63,0.69}{\textbf{\textit{{#1}}}}}
    \newcommand{\WarningTok}[1]{\textcolor[rgb]{0.38,0.63,0.69}{\textbf{\textit{{#1}}}}}
    \makeatletter
    \newsavebox\pandoc@box
    \newcommand*\pandocbounded[1]{%
      \sbox\pandoc@box{#1}%
      % scaling factors for width and height
      \Gscale@div\@tempa\textheight{\dimexpr\ht\pandoc@box+\dp\pandoc@box\relax}%
      \Gscale@div\@tempb\linewidth{\wd\pandoc@box}%
      % select the smaller of both
      \ifdim\@tempb\p@<\@tempa\p@
        \let\@tempa\@tempb
      \fi
      % scaling accordingly (\@tempa < 1)
      \ifdim\@tempa\p@<\p@
        \scalebox{\@tempa}{\usebox\pandoc@box}%
      % scaling not needed, use as it is
      \else
        \usebox{\pandoc@box}%
      \fi
    }
    \makeatother

    % Define a nice break command that doesn't care if a line doesn't already
    % exist.
    \def\br{\hspace*{\fill} \\* }
    % Math Jax compatibility definitions
    \def\gt{>}
    \def\lt{<}
    \let\Oldtex\TeX
    \let\Oldlatex\LaTeX
    \renewcommand{\TeX}{\textrm{\Oldtex}}
    \renewcommand{\LaTeX}{\textrm{\Oldlatex}}
    % Document parameters
    % Document title
    \title{Hughes\_assignment01}
    \author{by Ryan Hughes\\ Northeastern University\\ Master's in Complex Network Analysis}
    \date{\today}
    
    
    
    
    
    
% Pygments definitions
\makeatletter
\def\PY@reset{\let\PY@it=\relax \let\PY@bf=\relax%
    \let\PY@ul=\relax \let\PY@tc=\relax%
    \let\PY@bc=\relax \let\PY@ff=\relax}
\def\PY@tok#1{\csname PY@tok@#1\endcsname}
\def\PY@toks#1+{\ifx\relax#1\empty\else%
    \PY@tok{#1}\expandafter\PY@toks\fi}
\def\PY@do#1{\PY@bc{\PY@tc{\PY@ul{%
    \PY@it{\PY@bf{\PY@ff{#1}}}}}}}
\def\PY#1#2{\PY@reset\PY@toks#1+\relax+\PY@do{#2}}

\@namedef{PY@tok@w}{\def\PY@tc##1{\textcolor[rgb]{0.73,0.73,0.73}{##1}}}
\@namedef{PY@tok@c}{\let\PY@it=\textit\def\PY@tc##1{\textcolor[rgb]{0.24,0.48,0.48}{##1}}}
\@namedef{PY@tok@cp}{\def\PY@tc##1{\textcolor[rgb]{0.61,0.40,0.00}{##1}}}
\@namedef{PY@tok@k}{\let\PY@bf=\textbf\def\PY@tc##1{\textcolor[rgb]{0.00,0.50,0.00}{##1}}}
\@namedef{PY@tok@kp}{\def\PY@tc##1{\textcolor[rgb]{0.00,0.50,0.00}{##1}}}
\@namedef{PY@tok@kt}{\def\PY@tc##1{\textcolor[rgb]{0.69,0.00,0.25}{##1}}}
\@namedef{PY@tok@o}{\def\PY@tc##1{\textcolor[rgb]{0.40,0.40,0.40}{##1}}}
\@namedef{PY@tok@ow}{\let\PY@bf=\textbf\def\PY@tc##1{\textcolor[rgb]{0.67,0.13,1.00}{##1}}}
\@namedef{PY@tok@nb}{\def\PY@tc##1{\textcolor[rgb]{0.00,0.50,0.00}{##1}}}
\@namedef{PY@tok@nf}{\def\PY@tc##1{\textcolor[rgb]{0.00,0.00,1.00}{##1}}}
\@namedef{PY@tok@nc}{\let\PY@bf=\textbf\def\PY@tc##1{\textcolor[rgb]{0.00,0.00,1.00}{##1}}}
\@namedef{PY@tok@nn}{\let\PY@bf=\textbf\def\PY@tc##1{\textcolor[rgb]{0.00,0.00,1.00}{##1}}}
\@namedef{PY@tok@ne}{\let\PY@bf=\textbf\def\PY@tc##1{\textcolor[rgb]{0.80,0.25,0.22}{##1}}}
\@namedef{PY@tok@nv}{\def\PY@tc##1{\textcolor[rgb]{0.10,0.09,0.49}{##1}}}
\@namedef{PY@tok@no}{\def\PY@tc##1{\textcolor[rgb]{0.53,0.00,0.00}{##1}}}
\@namedef{PY@tok@nl}{\def\PY@tc##1{\textcolor[rgb]{0.46,0.46,0.00}{##1}}}
\@namedef{PY@tok@ni}{\let\PY@bf=\textbf\def\PY@tc##1{\textcolor[rgb]{0.44,0.44,0.44}{##1}}}
\@namedef{PY@tok@na}{\def\PY@tc##1{\textcolor[rgb]{0.41,0.47,0.13}{##1}}}
\@namedef{PY@tok@nt}{\let\PY@bf=\textbf\def\PY@tc##1{\textcolor[rgb]{0.00,0.50,0.00}{##1}}}
\@namedef{PY@tok@nd}{\def\PY@tc##1{\textcolor[rgb]{0.67,0.13,1.00}{##1}}}
\@namedef{PY@tok@s}{\def\PY@tc##1{\textcolor[rgb]{0.73,0.13,0.13}{##1}}}
\@namedef{PY@tok@sd}{\let\PY@it=\textit\def\PY@tc##1{\textcolor[rgb]{0.73,0.13,0.13}{##1}}}
\@namedef{PY@tok@si}{\let\PY@bf=\textbf\def\PY@tc##1{\textcolor[rgb]{0.64,0.35,0.47}{##1}}}
\@namedef{PY@tok@se}{\let\PY@bf=\textbf\def\PY@tc##1{\textcolor[rgb]{0.67,0.36,0.12}{##1}}}
\@namedef{PY@tok@sr}{\def\PY@tc##1{\textcolor[rgb]{0.64,0.35,0.47}{##1}}}
\@namedef{PY@tok@ss}{\def\PY@tc##1{\textcolor[rgb]{0.10,0.09,0.49}{##1}}}
\@namedef{PY@tok@sx}{\def\PY@tc##1{\textcolor[rgb]{0.00,0.50,0.00}{##1}}}
\@namedef{PY@tok@m}{\def\PY@tc##1{\textcolor[rgb]{0.40,0.40,0.40}{##1}}}
\@namedef{PY@tok@gh}{\let\PY@bf=\textbf\def\PY@tc##1{\textcolor[rgb]{0.00,0.00,0.50}{##1}}}
\@namedef{PY@tok@gu}{\let\PY@bf=\textbf\def\PY@tc##1{\textcolor[rgb]{0.50,0.00,0.50}{##1}}}
\@namedef{PY@tok@gd}{\def\PY@tc##1{\textcolor[rgb]{0.63,0.00,0.00}{##1}}}
\@namedef{PY@tok@gi}{\def\PY@tc##1{\textcolor[rgb]{0.00,0.52,0.00}{##1}}}
\@namedef{PY@tok@gr}{\def\PY@tc##1{\textcolor[rgb]{0.89,0.00,0.00}{##1}}}
\@namedef{PY@tok@ge}{\let\PY@it=\textit}
\@namedef{PY@tok@gs}{\let\PY@bf=\textbf}
\@namedef{PY@tok@ges}{\let\PY@bf=\textbf\let\PY@it=\textit}
\@namedef{PY@tok@gp}{\let\PY@bf=\textbf\def\PY@tc##1{\textcolor[rgb]{0.00,0.00,0.50}{##1}}}
\@namedef{PY@tok@go}{\def\PY@tc##1{\textcolor[rgb]{0.44,0.44,0.44}{##1}}}
\@namedef{PY@tok@gt}{\def\PY@tc##1{\textcolor[rgb]{0.00,0.27,0.87}{##1}}}
\@namedef{PY@tok@err}{\def\PY@bc##1{{\setlength{\fboxsep}{\string -\fboxrule}\fcolorbox[rgb]{1.00,0.00,0.00}{1,1,1}{\strut ##1}}}}
\@namedef{PY@tok@kc}{\let\PY@bf=\textbf\def\PY@tc##1{\textcolor[rgb]{0.00,0.50,0.00}{##1}}}
\@namedef{PY@tok@kd}{\let\PY@bf=\textbf\def\PY@tc##1{\textcolor[rgb]{0.00,0.50,0.00}{##1}}}
\@namedef{PY@tok@kn}{\let\PY@bf=\textbf\def\PY@tc##1{\textcolor[rgb]{0.00,0.50,0.00}{##1}}}
\@namedef{PY@tok@kr}{\let\PY@bf=\textbf\def\PY@tc##1{\textcolor[rgb]{0.00,0.50,0.00}{##1}}}
\@namedef{PY@tok@bp}{\def\PY@tc##1{\textcolor[rgb]{0.00,0.50,0.00}{##1}}}
\@namedef{PY@tok@fm}{\def\PY@tc##1{\textcolor[rgb]{0.00,0.00,1.00}{##1}}}
\@namedef{PY@tok@vc}{\def\PY@tc##1{\textcolor[rgb]{0.10,0.09,0.49}{##1}}}
\@namedef{PY@tok@vg}{\def\PY@tc##1{\textcolor[rgb]{0.10,0.09,0.49}{##1}}}
\@namedef{PY@tok@vi}{\def\PY@tc##1{\textcolor[rgb]{0.10,0.09,0.49}{##1}}}
\@namedef{PY@tok@vm}{\def\PY@tc##1{\textcolor[rgb]{0.10,0.09,0.49}{##1}}}
\@namedef{PY@tok@sa}{\def\PY@tc##1{\textcolor[rgb]{0.73,0.13,0.13}{##1}}}
\@namedef{PY@tok@sb}{\def\PY@tc##1{\textcolor[rgb]{0.73,0.13,0.13}{##1}}}
\@namedef{PY@tok@sc}{\def\PY@tc##1{\textcolor[rgb]{0.73,0.13,0.13}{##1}}}
\@namedef{PY@tok@dl}{\def\PY@tc##1{\textcolor[rgb]{0.73,0.13,0.13}{##1}}}
\@namedef{PY@tok@s2}{\def\PY@tc##1{\textcolor[rgb]{0.73,0.13,0.13}{##1}}}
\@namedef{PY@tok@sh}{\def\PY@tc##1{\textcolor[rgb]{0.73,0.13,0.13}{##1}}}
\@namedef{PY@tok@s1}{\def\PY@tc##1{\textcolor[rgb]{0.73,0.13,0.13}{##1}}}
\@namedef{PY@tok@mb}{\def\PY@tc##1{\textcolor[rgb]{0.40,0.40,0.40}{##1}}}
\@namedef{PY@tok@mf}{\def\PY@tc##1{\textcolor[rgb]{0.40,0.40,0.40}{##1}}}
\@namedef{PY@tok@mh}{\def\PY@tc##1{\textcolor[rgb]{0.40,0.40,0.40}{##1}}}
\@namedef{PY@tok@mi}{\def\PY@tc##1{\textcolor[rgb]{0.40,0.40,0.40}{##1}}}
\@namedef{PY@tok@il}{\def\PY@tc##1{\textcolor[rgb]{0.40,0.40,0.40}{##1}}}
\@namedef{PY@tok@mo}{\def\PY@tc##1{\textcolor[rgb]{0.40,0.40,0.40}{##1}}}
\@namedef{PY@tok@ch}{\let\PY@it=\textit\def\PY@tc##1{\textcolor[rgb]{0.24,0.48,0.48}{##1}}}
\@namedef{PY@tok@cm}{\let\PY@it=\textit\def\PY@tc##1{\textcolor[rgb]{0.24,0.48,0.48}{##1}}}
\@namedef{PY@tok@cpf}{\let\PY@it=\textit\def\PY@tc##1{\textcolor[rgb]{0.24,0.48,0.48}{##1}}}
\@namedef{PY@tok@c1}{\let\PY@it=\textit\def\PY@tc##1{\textcolor[rgb]{0.24,0.48,0.48}{##1}}}
\@namedef{PY@tok@cs}{\let\PY@it=\textit\def\PY@tc##1{\textcolor[rgb]{0.24,0.48,0.48}{##1}}}

\def\PYZbs{\char`\\}
\def\PYZus{\char`\_}
\def\PYZob{\char`\{}
\def\PYZcb{\char`\}}
\def\PYZca{\char`\^}
\def\PYZam{\char`\&}
\def\PYZlt{\char`\<}
\def\PYZgt{\char`\>}
\def\PYZsh{\char`\#}
\def\PYZpc{\char`\%}
\def\PYZdl{\char`\$}
\def\PYZhy{\char`\-}
\def\PYZsq{\char`\'}
\def\PYZdq{\char`\"}
\def\PYZti{\char`\~}
% for compatibility with earlier versions
\def\PYZat{@}
\def\PYZlb{[}
\def\PYZrb{]}
\makeatother


    % For linebreaks inside Verbatim environment from package fancyvrb.
    \makeatletter
        \newbox\Wrappedcontinuationbox
        \newbox\Wrappedvisiblespacebox
        \newcommand*\Wrappedvisiblespace {\textcolor{red}{\textvisiblespace}}
        \newcommand*\Wrappedcontinuationsymbol {\textcolor{red}{\llap{\tiny$\m@th\hookrightarrow$}}}
        \newcommand*\Wrappedcontinuationindent {3ex }
        \newcommand*\Wrappedafterbreak {\kern\Wrappedcontinuationindent\copy\Wrappedcontinuationbox}
        % Take advantage of the already applied Pygments mark-up to insert
        % potential linebreaks for TeX processing.
        %        {, <, #, %, $, ' and ": go to next line.
        %        _, }, ^, &, >, - and ~: stay at end of broken line.
        % Use of \textquotesingle for straight quote.
        \newcommand*\Wrappedbreaksatspecials {%
            \def\PYGZus{\discretionary{\char`\_}{\Wrappedafterbreak}{\char`\_}}%
            \def\PYGZob{\discretionary{}{\Wrappedafterbreak\char`\{}{\char`\{}}%
            \def\PYGZcb{\discretionary{\char`\}}{\Wrappedafterbreak}{\char`\}}}%
            \def\PYGZca{\discretionary{\char`\^}{\Wrappedafterbreak}{\char`\^}}%
            \def\PYGZam{\discretionary{\char`\&}{\Wrappedafterbreak}{\char`\&}}%
            \def\PYGZlt{\discretionary{}{\Wrappedafterbreak\char`\<}{\char`\<}}%
            \def\PYGZgt{\discretionary{\char`\>}{\Wrappedafterbreak}{\char`\>}}%
            \def\PYGZsh{\discretionary{}{\Wrappedafterbreak\char`\#}{\char`\#}}%
            \def\PYGZpc{\discretionary{}{\Wrappedafterbreak\char`\%}{\char`\%}}%
            \def\PYGZdl{\discretionary{}{\Wrappedafterbreak\char`\$}{\char`\$}}%
            \def\PYGZhy{\discretionary{\char`\-}{\Wrappedafterbreak}{\char`\-}}%
            \def\PYGZsq{\discretionary{}{\Wrappedafterbreak\textquotesingle}{\textquotesingle}}%
            \def\PYGZdq{\discretionary{}{\Wrappedafterbreak\char`\"}{\char`\"}}%
            \def\PYGZti{\discretionary{\char`\~}{\Wrappedafterbreak}{\char`\~}}%
        }
        % Some characters . , ; ? ! / are not pygmentized.
        % This macro makes them "active" and they will insert potential linebreaks
        \newcommand*\Wrappedbreaksatpunct {%
            \lccode`\~`\.\lowercase{\def~}{\discretionary{\hbox{\char`\.}}{\Wrappedafterbreak}{\hbox{\char`\.}}}%
            \lccode`\~`\,\lowercase{\def~}{\discretionary{\hbox{\char`\,}}{\Wrappedafterbreak}{\hbox{\char`\,}}}%
            \lccode`\~`\;\lowercase{\def~}{\discretionary{\hbox{\char`\;}}{\Wrappedafterbreak}{\hbox{\char`\;}}}%
            \lccode`\~`\:\lowercase{\def~}{\discretionary{\hbox{\char`\:}}{\Wrappedafterbreak}{\hbox{\char`\:}}}%
            \lccode`\~`\?\lowercase{\def~}{\discretionary{\hbox{\char`\?}}{\Wrappedafterbreak}{\hbox{\char`\?}}}%
            \lccode`\~`\!\lowercase{\def~}{\discretionary{\hbox{\char`\!}}{\Wrappedafterbreak}{\hbox{\char`\!}}}%
            \lccode`\~`\/\lowercase{\def~}{\discretionary{\hbox{\char`\/}}{\Wrappedafterbreak}{\hbox{\char`\/}}}%
            \catcode`\.\active
            \catcode`\,\active
            \catcode`\;\active
            \catcode`\:\active
            \catcode`\?\active
            \catcode`\!\active
            \catcode`\/\active
            \lccode`\~`\~
        }
    \makeatother

    \let\OriginalVerbatim=\Verbatim
    \makeatletter
    \renewcommand{\Verbatim}[1][1]{%
        %\parskip\z@skip
        \sbox\Wrappedcontinuationbox {\Wrappedcontinuationsymbol}%
        \sbox\Wrappedvisiblespacebox {\FV@SetupFont\Wrappedvisiblespace}%
        \def\FancyVerbFormatLine ##1{\hsize\linewidth
            \vtop{\raggedright\hyphenpenalty\z@\exhyphenpenalty\z@
                \doublehyphendemerits\z@\finalhyphendemerits\z@
                \strut ##1\strut}%
        }%
        % If the linebreak is at a space, the latter will be displayed as visible
        % space at end of first line, and a continuation symbol starts next line.
        % Stretch/shrink are however usually zero for typewriter font.
        \def\FV@Space {%
            \nobreak\hskip\z@ plus\fontdimen3\font minus\fontdimen4\font
            \discretionary{\copy\Wrappedvisiblespacebox}{\Wrappedafterbreak}
            {\kern\fontdimen2\font}%
        }%

        % Allow breaks at special characters using \PYG... macros.
        \Wrappedbreaksatspecials
        % Breaks at punctuation characters . , ; ? ! and / need catcode=\active
        \OriginalVerbatim[#1,codes*=\Wrappedbreaksatpunct]%
    }
    \makeatother

    % Exact colors from NB
    \definecolor{incolor}{HTML}{303F9F}
    \definecolor{outcolor}{HTML}{D84315}
    \definecolor{cellborder}{HTML}{CFCFCF}
    \definecolor{cellbackground}{HTML}{F7F7F7}

    % prompt
    \makeatletter
    \newcommand{\boxspacing}{\kern\kvtcb@left@rule\kern\kvtcb@boxsep}
    \makeatother
    \newcommand{\prompt}[4]{
        {\ttfamily\llap{{\color{#2}[#3]:\hspace{3pt}#4}}\vspace{-\baselineskip}}
    }
    

    
    % Prevent overflowing lines due to hard-to-break entities
    \sloppy
    % Setup hyperref package
    \hypersetup{
      breaklinks=true,  % so long urls are correctly broken across lines
      colorlinks=true,
      urlcolor=urlcolor,
      linkcolor=linkcolor,
      citecolor=citecolor,
      }
    % Slightly bigger margins than the latex defaults
    
    \geometry{verbose,tmargin=1in,bmargin=1in,lmargin=1in,rmargin=1in}
    
    

\begin{document}
    
    \maketitle
    
    

    
    \paragraph{Question 0 (0 points)}\label{question-0-0-points}

Include your code in a (well-documented, commented, clear)
assignment01.ipynb file. Scores for the code-based questions below will
be based on your written responses and code

    \paragraph{Question 1 (20 points)}\label{question-1-20-points}

There are a lot of ways to store network data. Each format generally has
advantages and disadvantages that make them suited to particular types
of situations. Let's explore the differences between storing network
data in a .csv, pickle, or .gml file, and in Elasticsearch, which is
commonly used for graph database applications.

\emph{For your research, look at the documentation for each approach,
seek out StackOverflow threads, and consider consulting the Bagrow \&
Ahn textbook.}

    \begin{enumerate}
\def\labelenumi{(\alph{enumi})}
\tightlist
\item
  What are the advantages of each storage approach?
\end{enumerate}

    \begin{enumerate}
\def\labelenumi{\arabic{enumi}.}
\tightlist
\item
  .csv according to
  https://stackoverflow.com/questions/48770542/what-is-the-difference-between-save-a-pandas-dataframe-to-pickle-and-to-csv,
  has the advantage of being human readable and cross platform
\item
  pickle according to
  https://stackoverflow.com/questions/48770542/what-is-the-difference-between-save-a-pandas-dataframe-to-pickle-and-to-csv
  has the advantage of fast saving and loading and being more space
  efficient (for disk saves) than csv. It also saves the exact
  representation of the dataframe per how the pandas library sees it,
  thereby reducing the risk of lost data when converting to csv and
  back, and can store any binary subject.
\item
  .gml, according to Gephi
  (https://docs.gephi.org/desktop/User\_Manual/Import/GML\_File\_Format/)
  has the advantage of being a plain text file format with a simple,
  human-readable syntax used to represent graphs. Storing all nodes,
  edges and attributes into a single file
\item
  Elastisearch
  (https://stackoverflow.com/questions/57328151/how-does-elasticsearch-store-data)
  is a vector database solution. It stores data via ``lucene'' a text
  search engine, that saves data in a custom binary format. It has the
  advantage of being highly optimized. It also can handle any kind of
  data: structured, unstructured and vectors, and allows for search
  engine, analytics engine, geospatial engine and inference service
  queries on your dataset.
\end{enumerate}

    \begin{enumerate}
\def\labelenumi{(\alph{enumi})}
\setcounter{enumi}{1}
\tightlist
\item
  What are the disadvantages of each storage approach?
\end{enumerate}

\begin{enumerate}
\def\labelenumi{\arabic{enumi}.}
\tightlist
\item
  .csv according to the textbook can be subject to data loss due to the
  conversion process. In particular, delimiter collisions where part of
  the data contains the delimiter innately (e.g., ``,'' in a ``Surname,
  First Name'') used to separate entries. In addition, .csv is less
  space efficient than other formats.
\item
  pickle according to
  https://stackoverflow.com/questions/48770542/what-is-the-difference-between-save-a-pandas-dataframe-to-pickle-and-to-csv
  has the disadvantage of not being as universal and cross-platform as
  other formats, it may require you to be using the same Python version
  to save between computers for instance, and is generally not binary
  stable. It also can introduce potential security vulnerabilities
  should you be downloading an unknown file in this format that was
  corrupted or contains malicious code,
\item
  .gml according to the textbook has the disadvantage of using a
  non-standard format. It also is not fully supported by Gephi according
  to their documentation (particularly for some data structures
  \emph{though it's not specified which}). It also has the disadvantage
  of being similarly named to GraphML and Geography Markup Language
  (also GML).
\item
  Elastisearch, according to GeeksForGeeks
  (https://www.geeksforgeeks.org/elasticsearch/advantages-and-disadvantages-of-elasticsearch/)
  has the disadvantage of being very complex to setup and maintain as an
  industry grade database solution. It also is resource intensive
  requiring sufficient CPU, memory and storage. It also prioritizes
  performance and scale over consistency and reliability, meaning that
  there are scenarios where data is not 100\% reliable across the
  database machines. It also requires expertise and has a steep learning
  curve to pick-up.
\end{enumerate}

    \begin{enumerate}
\def\labelenumi{(\alph{enumi})}
\setcounter{enumi}{2}
\tightlist
\item
  At least one of these storage approaches poses a non-trivial security
  risk. Which one(s) should we worry about, and what is the security
  risk? Please explain this as simply as possible.
\end{enumerate}

    .pkl (pickle) may have a security problem according to the official
python documentation. The exact warning can be found here
https://docs.python.org/3/library/pickle.html as:

``\emph{Warning: The pickle module is not secure. Only unpickle data you
trust.}

\emph{It is possible to construct malicious pickle data which will
execute arbitrary code during unpickling. Never unpickle data that could
have come from an untrusted source, or that could have been tampered
with.}

\emph{Consider signing data with hmac if you need to ensure that it has
not been tampered with.}

\emph{Safer serialization formats such as json may be more appropriate
if you are processing untrusted data. See Comparison with json.}''

What this means is that when loading in a .pkl file, should that file
secretly contain code, Python may execute that code rather than simply
load it as a plain data-frame. As such, if the code is malicious, you
could unwittingly introduce hackers or malware to your pc.

    \begin{enumerate}
\def\labelenumi{(\alph{enumi})}
\setcounter{enumi}{3}
\tightlist
\item
  For each storage approach, think of a scenario where that storage
  approach would be an optimal choice. Explain your reasoning.
\end{enumerate}

    \begin{enumerate}
\def\labelenumi{\arabic{enumi}.}
\tightlist
\item
  .csv can be good when you are working with simple data that you know
  for certain would not contain potential collisions (e.g., data that is
  all integers perhaps). It would also be optimal in a scenario when you
  need to share data widely to many different machines and programs.
\item
  pickle can be optimal when you are working with your own data and need
  a fast and reliable way to store it. If you create the file and do not
  share it to potential bad actors, you can avoid the security risk
  related to this format, and thereby take advantage of its reliability,
  speed and spatial efficiency.
\item
  .gml can be optimal when you don't need to share the data very widely
  but need some of the benefits of .csv (e.g., human-readable) and,
  cannot store it as such due to delimiter collisions. .gml provides a
  similar human-readable benefit with the tradeoff of less universality.
\item
  Elastisearch, can be optimal when you truly need a long-term
  industry-grade solution for a large dataset, and want the performance,
  scalability and unique analytics engine capabilities that can come
  from this.
\end{enumerate}

    \paragraph{Question 2 (20 points).}\label{question-2-20-points.}

\emph{Algorithm 11.1 in Bagrow \& Ahn}, the basic edge exchange
algorithm, modifies a graph in place to produce a null model while
preserving the degree distribution. It does so by picking two edges at
random and swapping their endpoints.

    \begin{itemize}
\tightlist
\item
  1: procedure EDGE EXCHANGE(G, t)
\item
  2: s ← 0
\item
  3: while s \textless{} t do
\item
  4: (i, j) ← get random edge(G)
\item
  2
\item
  5: (u, v) ← get random edge(G)
\item
  6: if i = u or i = v or j = u or j = v then
\item
  7: continue
\item
  8: end if
\item
  9: if has edge(G, i, v) or has edge(G, j, u) then
\item
  10: continue
\item
  11: end if
\item
  12: remove edge(G, i, j)
\item
  13: remove edge(G, u, v)
\item
  14: add edge(G, i, v)
\item
  15: add edge(G, j, u)
\item
  16: s ← s + 1
\item
  17: end while
\item
  18: return G
\item
  19: end procedure
\end{itemize}

    \begin{enumerate}
\def\labelenumi{(\alph{enumi})}
\tightlist
\item
  (Bagrow \& Ahn Question 11.2) The basic edge exchange algorithm (as
  shown above) will enter an infinite loop and never end if no allowable
  exchanges are present. Describe a scenario where you will never be
  able to find an allowable edge exchange according to this algorithm.
\end{enumerate}

    One simple edge-case scenario when you will never be able to find an
allowable edge exchange according to this algorithm is, if you're
dataset is too small. For instance, in theory, if you had a graph with
just 1 edge between 2 nodes, this algorithm may never find a pair to
exchange,

    \begin{enumerate}
\def\labelenumi{(\alph{enumi})}
\setcounter{enumi}{1}
\tightlist
\item
  (Bagrow \& Ahn Question 11.3) Modify the algorithm to prevent an
  infinite loop by abandoning the exchange process after reaching a
  given number of attempts without a successful exchange.
\end{enumerate}

    \begin{itemize}
\tightlist
\item
  1: procedure EDGE EXCHANGE(G, t) \#Hard time-out when we have missed a
  total number of times equal to the sum of all edges in the graph.
\item
  \textbf{1b: hard\_limit = get\_total\_number\_of\_edges(G)}
\item
  2: s ← 0
\item
  3: while s \textless{} t and hard\_limit \textgreater{} 0 do
\item
  4: (i, j) ← get random edge(G)
\item
  5: (u, v) ← get random edge(G)
\item
  6: if i = u or i = v or j = u or j = v then
\item
  6b: \textbf{hard\_limit -=1}
\item
  7: continue
\item
  8: end if
\item
  9: if has edge(G, i, v) or has edge(G, j, u) then
\item
  9b: \textbf{hard\_limit -=1}
\item
  10: continue
\item
  11: end if
\item
  12: remove edge(G, i, j)
\item
  13: remove edge(G, u, v)
\item
  14: add edge(G, i, v)
\item
  15: add edge(G, j, u)
\item
  16: s ← s + 1
\item
  17: end while
\item
  18: return G
\item
  19: end procedure
\end{itemize}

    \begin{enumerate}
\def\labelenumi{(\alph{enumi})}
\setcounter{enumi}{2}
\tightlist
\item
  (Bagrow \& Ahn Question 11.4) Modify the algorithm to preserve
  \textbf{not only the degrees of every node but also the overall
  connectivity of the network}. How does this additional constraint make
  edge exchange more difficult?
\end{enumerate}

    \begin{itemize}
\tightlist
\item
  1: procedure EDGE EXCHANGE(G, t)
\item
  1b: hard\_limit = get\_total\_number\_of\_edges(G) \# Hard time-out
  when we have missed a total number of times equal to the sum of all
  edges in the graph.
\item
  \textbf{1c: components = nx.connected\_components(G)}
\item
  2: s ← 0
\item
  3: while s \textless{} t and hard\_limit \textgreater{} 0 do
\item
  4: (i, j) ← get random edge(G)
\item
  5: (u, v) ← get random edge(G)
\item
  5b: if not unique(i,j,u,v) do
\item
  5c: (i, j) ← get random edge(G)
\item
  5d: (u, v) ← get random edge(G)
\item
  5e: hard\_limit -=1
\item
  5f: end if
\item
  6: if i = u or i = v or j = u or j = v then
\item
  6b: hard\_limit -=1
\item
  7: continue
\item
  8: end if
\item
  9: if has edge(G, i, v) or has edge(G, j, u) then
\item
  9b: hard\_limit -=1
\item
  10: continue
\item
  11: end if
\item
  11a: all\_in\_the\_same = True
\item
  11b: \textbf{for component in components:}
\item
  11c: \textbf{if(v in component and i in component) and (u in component
  and j in component)) or (v not in component and i not in component)
  and (u not in component and j not in component)):}
\item
  11d: \textbf{all\_in\_the\_same = True}
\item
  11e: \textbf{end if}
\item
  11f: \textbf{else:}
\item
  11g: \textbf{all\_in\_the\_same = False}
\item
  11h: \textbf{hard\_limit -=1}
\item
  11i: \textbf{continue}
\item
  11j: \textbf{end else}
\item
  11k: \textbf{end for}
\item
  11m: \textbf{if all\_in\_the\_same:} \#This ensures prior connectivity
  remains by checking that all nodes are all apart, or all not apart, of
  each component
\item
  12: remove edge(G, i, j)
\item
  13: remove edge(G, u, v)
\item
  14: add edge(G, i, v)
\item
  15: add edge(G, j, u)
\item
  15b: \textbf{end if}
\item
  16: s ← s + 1
\item
  17: end while
\item
  18: return G
\item
  19: end procedure
\end{itemize}

\emph{This additional constraint makes the edge constraint much more
difficult because it requires that all 4 nodes (corresponding to each
end of each pair) in addition to must being unique, in order to prevent
self-loops or parallel edges, must also be apart of all the same
components in the graph, in order to maintain connectivity levels.}

    \begin{enumerate}
\def\labelenumi{(\alph{enumi})}
\setcounter{enumi}{3}
\tightlist
\item
  (Bagrow \& Ahn Question 11.5) Suppose your network has a binary
  attribute xi ∈ \{−1, 1\} associated with every node i. Modify the
  algorithm to preserve the correlation on edges of x. This means that
  any new edge (a, b) inserted in place of an existing edge (i, j) must
  have the same attributes (xi, xj ) = (xa, xb). How might this
  additional constraint make edge exchange more difficult?
\end{enumerate}

    \begin{itemize}
\item
  1: procedure EDGE EXCHANGE(G, t)
\item
  1b: hard\_limit = get\_total\_number\_of\_edges(G) \# Hard time-out
  when we have missed a total number of times equal to the sum of all
  edges in the graph.
\item
  \textbf{1c: attribute = nx.get\_node\_attributes(G)}
\item
  2: s ← 0
\item
  3: while s \textless{} t and hard\_limit \textgreater{} 0 do
\item
  4: (i, j) ← get random edge(G)
\item
  5: (u, v) ← get random edge(G)
\item
  5b: if not unique(i,j,u,v) do
\item
  5c: (i, j) ← get random edge(G)
\item
  5d: (u, v) ← get random edge(G)
\item
  5e: hard\_limit -=1
\item
  5f: end if
\item
  \textbf{5g: if ( attribute{[}i{]}!=attribute{[}u{]} or
  attribute{[}i{]}!=attribute{[}v{]} or}
  \textbf{attribute{[}j{]}!=attribute{[}u{]} or
  attribute{[}j{]}!=attribute{[}v{]} )}
\item
  **5h: continue
\item
  \textbf{5i: hard\_limit -=1}
\item
  \textbf{5j: end if}\\
\item
  6: if i = u or i = v or j = u or j = v then
\item
  6b: hard\_limit -=1
\item
  7: continue
\item
  8: end if
\item
  9: if has edge(G, i, v) or has edge(G, j, u) then
\item
  9b: hard\_limit -=1
\item
  10: continue
\item
  11: end if
\item
  12: remove edge(G, i, j)
\item
  13: remove edge(G, u, v)
\item
  14: add edge(G, i, v)
\item
  15: add edge(G, j, u)
\item
  16: s ← s + 1
\item
  17: end while
\item
  18: return G
\item
  19: end procedure
\end{itemize}

\emph{This additional constraint makes the edge constraint much more
difficult in a similar way to the above algorithm for maintaining
connectivity. It requires that the new edges to be made consist of the
same attribute pairs as the previous attribute pairs. Because we are
removing a specific pair of attributes as a part of replacement, the
effect is that all attrbitues must be equivalent in order to preserve
the correlation on edges of x. Depending on how rare this attribute
status is the number of swaps possible may be small. Note: This logic
could be cleaned up a bit into if-else statements to make each check
prior to each possible swap rather than all at once.}

    \subsubsection{\texorpdfstring{\textbf{Question 3 (20
points).}}{Question 3 (20 points).}}\label{question-3-20-points.}

There are many ways to measure a node's centrality in a network. Pick
one centrality measure out of \emph{\{betweenness centrality, harmonic
centrality, percolation centrality\}} and do the following:

    \begin{enumerate}
\def\labelenumi{(\alph{enumi})}
\tightlist
\item
  Write your own function to compute a single node's centrality in a
  graph. This function should take a networkx graph and one node in the
  graph as input. As output, it should return a float or integer
  representing the node's centrality in the graph according to one of
  the centrality measures listed above.
\end{enumerate}

    \begin{tcolorbox}[breakable, size=fbox, boxrule=1pt, pad at break*=1mm,colback=cellbackground, colframe=cellborder]
\prompt{In}{incolor}{4}{\boxspacing}
\begin{Verbatim}[commandchars=\\\{\}]
\PY{k+kn}{import}\PY{+w}{ }\PY{n+nn}{networkx}\PY{+w}{ }\PY{k}{as}\PY{+w}{ }\PY{n+nn}{nx}
\PY{k+kn}{from}\PY{+w}{ }\PY{n+nn}{collections}\PY{+w}{ }\PY{k+kn}{import} \PY{n}{Counter}

\PY{k}{def}\PY{+w}{ }\PY{n+nf}{configuration\PYZus{}model\PYZus{}from\PYZus{}degree\PYZus{}sequence}\PY{p}{(}\PY{n}{degree\PYZus{}sequence}\PY{p}{,} \PY{n}{return\PYZus{}simple}\PY{o}{=}\PY{k+kc}{True}\PY{p}{)}\PY{p}{:}
\PY{+w}{    }\PY{l+s+sd}{\PYZdq{}\PYZdq{}\PYZdq{}}
\PY{l+s+sd}{    Generate a random graph using the configuration model from a given degree sequence}
\PY{l+s+sd}{    without using the NetworkX built\PYZhy{}in function.}

\PY{l+s+sd}{    The configuration model generates a random graph that preserves the degree }
\PY{l+s+sd}{    sequence of nodes by assigning \PYZdq{}stubs\PYZdq{} or \PYZdq{}half\PYZhy{}edges\PYZdq{} to each node and }
\PY{l+s+sd}{    randomly pairing these stubs to form edges. This process can result in }
\PY{l+s+sd}{    graphs with self\PYZhy{}loops and parallel edges, which can be removed if needed.}

\PY{l+s+sd}{    Parameters}
\PY{l+s+sd}{    \PYZhy{}\PYZhy{}\PYZhy{}\PYZhy{}\PYZhy{}\PYZhy{}\PYZhy{}\PYZhy{}\PYZhy{}\PYZhy{}}
\PY{l+s+sd}{    degree\PYZus{}sequence : list of int}
\PY{l+s+sd}{        A list representing the degree of each node in the graph. The sum of }
\PY{l+s+sd}{        the degrees in this sequence must be even for the configuration model }
\PY{l+s+sd}{        to create a valid graph.}

\PY{l+s+sd}{    Returns}
\PY{l+s+sd}{    \PYZhy{}\PYZhy{}\PYZhy{}\PYZhy{}\PYZhy{}\PYZhy{}\PYZhy{}}
\PY{l+s+sd}{    G : networkx.MultiGraph}
\PY{l+s+sd}{        A multigraph generated from the given degree sequence. The graph may }
\PY{l+s+sd}{        contain self\PYZhy{}loops and parallel edges, which are allowed in the }
\PY{l+s+sd}{        configuration model.}

\PY{l+s+sd}{    Notes}
\PY{l+s+sd}{    \PYZhy{}\PYZhy{}\PYZhy{}\PYZhy{}\PYZhy{}}
\PY{l+s+sd}{    \PYZhy{} This method works by assigning \PYZdq{}stubs\PYZdq{} or \PYZdq{}half\PYZhy{}edges\PYZdq{} to nodes based on }
\PY{l+s+sd}{      their degree and then randomly pairing them to form edges. The resulting }
\PY{l+s+sd}{      graph can have self\PYZhy{}loops and parallel edges.}
\PY{l+s+sd}{    \PYZhy{} Self\PYZhy{}loops and parallel edges can be removed post\PYZhy{}generation if a simple }
\PY{l+s+sd}{      graph is required using NetworkX\PYZsq{}s `nx.Graph(G)`.}
\PY{l+s+sd}{    \PYZhy{} The degree sequence must have an even sum for a valid graph construction. }
\PY{l+s+sd}{      If the sum of the degrees is odd, no graph can be constructed.}

\PY{l+s+sd}{    Time Complexity}
\PY{l+s+sd}{    \PYZhy{}\PYZhy{}\PYZhy{}\PYZhy{}\PYZhy{}\PYZhy{}\PYZhy{}\PYZhy{}\PYZhy{}\PYZhy{}\PYZhy{}\PYZhy{}\PYZhy{}\PYZhy{}\PYZhy{}}
\PY{l+s+sd}{    The time complexity is O(E), where E is the number of edges in the graph.}

\PY{l+s+sd}{    Example}
\PY{l+s+sd}{    \PYZhy{}\PYZhy{}\PYZhy{}\PYZhy{}\PYZhy{}\PYZhy{}\PYZhy{}}
\PY{l+s+sd}{    \PYZgt{}\PYZgt{}\PYZgt{} degree\PYZus{}sequence = [3, 3, 2, 2, 1, 1]}
\PY{l+s+sd}{    \PYZgt{}\PYZgt{}\PYZgt{} G = configuration\PYZus{}model\PYZus{}from\PYZus{}degree\PYZus{}sequence(degree\PYZus{}sequence)}
\PY{l+s+sd}{    \PYZgt{}\PYZgt{}\PYZgt{} nx.is\PYZus{}graphical(degree\PYZus{}sequence)}
\PY{l+s+sd}{    True}
\PY{l+s+sd}{    \PYZdq{}\PYZdq{}\PYZdq{}}

    \PY{c+c1}{\PYZsh{} Check if the degree sequence is valid (sum of degrees must be even)}
    \PY{k}{if} \PY{n+nb}{sum}\PY{p}{(}\PY{n}{degree\PYZus{}sequence}\PY{p}{)} \PY{o}{\PYZpc{}} \PY{l+m+mi}{2} \PY{o}{!=} \PY{l+m+mi}{0}\PY{p}{:}
        \PY{k}{raise} \PY{n+ne}{ValueError}\PY{p}{(}\PY{l+s+s2}{\PYZdq{}}\PY{l+s+s2}{The sum of the degree sequence must be even.}\PY{l+s+s2}{\PYZdq{}}\PY{p}{)}

    \PY{c+c1}{\PYZsh{} Create stubs list: node i appears degree\PYZus{}sequence[i] times}
    \PY{n}{stubs} \PY{o}{=} \PY{p}{[}\PY{p}{]}
    \PY{k}{for} \PY{n}{node}\PY{p}{,} \PY{n}{degree} \PY{o+ow}{in} \PY{n+nb}{enumerate}\PY{p}{(}\PY{n}{degree\PYZus{}sequence}\PY{p}{)}\PY{p}{:}
        \PY{n}{stubs}\PY{o}{.}\PY{n}{extend}\PY{p}{(}\PY{p}{[}\PY{n}{node}\PY{p}{]} \PY{o}{*} \PY{n}{degree}\PY{p}{)}

    \PY{c+c1}{\PYZsh{} Shuffle stubs to randomize the pairing process}
    \PY{n}{np}\PY{o}{.}\PY{n}{random}\PY{o}{.}\PY{n}{shuffle}\PY{p}{(}\PY{n}{stubs}\PY{p}{)}

    \PY{c+c1}{\PYZsh{} Initialize an empty multigraph}
    \PY{n}{G} \PY{o}{=} \PY{n}{nx}\PY{o}{.}\PY{n}{MultiGraph}\PY{p}{(}\PY{p}{)}

    \PY{c+c1}{\PYZsh{} Add nodes to the graph}
    \PY{n}{G}\PY{o}{.}\PY{n}{add\PYZus{}nodes\PYZus{}from}\PY{p}{(}\PY{n+nb}{range}\PY{p}{(}\PY{n+nb}{len}\PY{p}{(}\PY{n}{degree\PYZus{}sequence}\PY{p}{)}\PY{p}{)}\PY{p}{)}

    \PY{c+c1}{\PYZsh{} Pair stubs to create edges}
    \PY{k}{while} \PY{n}{stubs}\PY{p}{:}
        \PY{n}{u} \PY{o}{=} \PY{n}{stubs}\PY{o}{.}\PY{n}{pop}\PY{p}{(}\PY{p}{)}
        \PY{n}{v} \PY{o}{=} \PY{n}{stubs}\PY{o}{.}\PY{n}{pop}\PY{p}{(}\PY{p}{)}

        \PY{c+c1}{\PYZsh{} Add the edge to the graph}
        \PY{n}{G}\PY{o}{.}\PY{n}{add\PYZus{}edge}\PY{p}{(}\PY{n}{u}\PY{p}{,} \PY{n}{v}\PY{p}{)}

    \PY{k}{if} \PY{n}{return\PYZus{}simple}\PY{p}{:}
        \PY{c+c1}{\PYZsh{} Convert the multigraph to a simple graph (remove parallel edges and self\PYZhy{}loops)}
        \PY{n}{G\PYZus{}simple} \PY{o}{=} \PY{n}{nx}\PY{o}{.}\PY{n}{Graph}\PY{p}{(}\PY{n}{G}\PY{p}{)}  \PY{c+c1}{\PYZsh{} This removes parallel edges and self\PYZhy{}loops by default}

        \PY{k}{return} \PY{n}{G\PYZus{}simple}

    \PY{k}{else}\PY{p}{:}
        \PY{k}{return} \PY{n}{G}
\end{Verbatim}
\end{tcolorbox}

    \begin{tcolorbox}[breakable, size=fbox, boxrule=1pt, pad at break*=1mm,colback=cellbackground, colframe=cellborder]
\prompt{In}{incolor}{5}{\boxspacing}
\begin{Verbatim}[commandchars=\\\{\}]
\PY{k}{def}\PY{+w}{ }\PY{n+nf}{betweenness\PYZus{}centrality\PYZus{}from\PYZus{}scratch}\PY{p}{(}\PY{n}{G}\PY{p}{,} \PY{n}{node\PYZus{}i}\PY{p}{)}\PY{p}{:}
\PY{+w}{    }\PY{l+s+sd}{\PYZdq{}\PYZdq{}\PYZdq{}}
\PY{l+s+sd}{    Calculate the betweenness centrality for each node in a graph from scratch.}

\PY{l+s+sd}{    Betweenness centrality is a measure of the extent to which a node lies on the shortest}
\PY{l+s+sd}{    paths between other nodes in the network. This function computes the betweenness}
\PY{l+s+sd}{    centrality for all nodes in the graph `G`.}

\PY{l+s+sd}{    Parameters}
\PY{l+s+sd}{    \PYZhy{}\PYZhy{}\PYZhy{}\PYZhy{}\PYZhy{}\PYZhy{}\PYZhy{}\PYZhy{}\PYZhy{}\PYZhy{}}
\PY{l+s+sd}{    G : networkx.Graph}
\PY{l+s+sd}{        The input graph on which the betweenness centrality is calculated. It can}
\PY{l+s+sd}{        be any type of graph (undirected, directed, etc.) supported by NetworkX.}

\PY{l+s+sd}{    node\PYZus{}i : networkx.Node}
\PY{l+s+sd}{        The input node on which the betweenness centrality is calculated. }
\PY{l+s+sd}{        }
\PY{l+s+sd}{    Returns}
\PY{l+s+sd}{    \PYZhy{}\PYZhy{}\PYZhy{}\PYZhy{}\PYZhy{}\PYZhy{}\PYZhy{}}
\PY{l+s+sd}{    centrality : float}
\PY{l+s+sd}{        A float where the value is the corresponding betweenness centrality score as a function of the total number of nodes N, the total number of geodesics for all pairs of nodes in  graph G, and the number of geodesics node\PYZus{}i is part of in said graph, G.}


\PY{l+s+sd}{    Time Complexity}
\PY{l+s+sd}{    \PYZhy{}\PYZhy{}\PYZhy{}\PYZhy{}\PYZhy{}\PYZhy{}\PYZhy{}\PYZhy{}\PYZhy{}\PYZhy{}\PYZhy{}\PYZhy{}\PYZhy{}\PYZhy{}\PYZhy{}}
\PY{l+s+sd}{    The overall time complexity is O(N\PYZca{}2) where N is the number of nodes, due to the use of}
\PY{l+s+sd}{    combinations(G.nodes(), 2) dominating, which provides all \PYZdq{}N choose 2\PYZdq{} pairs of nodes.}

\PY{l+s+sd}{    Notes}
\PY{l+s+sd}{    \PYZhy{}\PYZhy{}\PYZhy{}\PYZhy{}\PYZhy{}}
\PY{l+s+sd}{    If unweighted, networkx\PYZsq{}s all\PYZus{}shortest\PYZus{}paths() uses Breadth\PYZhy{}First Search which is O(N), if weighted though, nx.all\PYZus{}shortest\PYZus{}paths() defaults to}
\PY{l+s+sd}{    Dijkstra\PYZsq{}s which is O(E*N*log(N)), and if negative weights exist then it is O(N*E) due to Bellman\PYZhy{}Ford. All where N is the number of nodes, where E is}
\PY{l+s+sd}{    the number of edges.}
\PY{l+s+sd}{    https://networkx.org/documentation/stable/\PYZus{}modules/networkx/algorithms/shortest\PYZus{}paths/unweighted.html}

\PY{l+s+sd}{    Citations}
\PY{l+s+sd}{    \PYZhy{}\PYZhy{}\PYZhy{}\PYZhy{}\PYZhy{}\PYZhy{}\PYZhy{}\PYZhy{}\PYZhy{}}
\PY{l+s+sd}{    Freeman, Linton C. “A Set of Measures of Centrality Based on Betweenness.” Sociometry, vol. 40, no. 1, 1977, pp. 35–41. JSTOR, https://doi.org/10.2307/3033543. Accessed 6 Oct. 2026.}

\PY{l+s+sd}{    Example}
\PY{l+s+sd}{    \PYZhy{}\PYZhy{}\PYZhy{}\PYZhy{}\PYZhy{}\PYZhy{}\PYZhy{}}
\PY{l+s+sd}{    \PYZgt{}\PYZgt{}\PYZgt{} import networkx as nx}
\PY{l+s+sd}{    \PYZgt{}\PYZgt{}\PYZgt{} G = nx.karate\PYZus{}club\PYZus{}graph()}
\PY{l+s+sd}{    \PYZgt{}\PYZgt{}\PYZgt{} node\PYZus{}i = 3}
\PY{l+s+sd}{    \PYZgt{}\PYZgt{}\PYZgt{} betweenness\PYZus{}centrality\PYZus{}from\PYZus{}scratch(kcg, node\PYZus{}i)}
\PY{l+s+sd}{    (3, 0.011909271284271286)}
\PY{l+s+sd}{    \PYZdq{}\PYZdq{}\PYZdq{}}
    \PY{k+kn}{import}\PY{+w}{ }\PY{n+nn}{networkx}\PY{+w}{ }\PY{k}{as}\PY{+w}{ }\PY{n+nn}{nx}
    \PY{k+kn}{from}\PY{+w}{ }\PY{n+nn}{itertools}\PY{+w}{ }\PY{k+kn}{import} \PY{n}{combinations}
    \PY{n}{centrality} \PY{o}{=} \PY{l+m+mf}{0.0}
    \PY{n}{N} \PY{o}{=} \PY{n}{G}\PY{o}{.}\PY{n}{number\PYZus{}of\PYZus{}nodes}\PY{p}{(}\PY{p}{)}  \PY{c+c1}{\PYZsh{} Total number of nodes in the graph}

\PY{+w}{    }\PY{l+s+sd}{\PYZsq{}\PYZsq{}\PYZsq{}}
\PY{l+s+sd}{    Grabbing C\PYZus{}B(p\PYZus{}k) =\PYZgt{} C\PYZus{}B(node\PYZus{}i) = \PYZdq{}the overall partial betweenness of point p\PYZus{}k.\PYZdq{}}
\PY{l+s+sd}{    If there is only 1 geodesic (i.e., shortest path) C\PYZus{}B(p\PYZus{}k) is increased by 1.}
\PY{l+s+sd}{    Otherwise, C\PYZus{}B(p\PYZus{}k) is increased relative to the proportion of geodesics p\PYZus{}k is apart of}
\PY{l+s+sd}{    \PYZsq{}\PYZsq{}\PYZsq{}}
    \PY{k}{for} \PY{n}{a}\PY{p}{,} \PY{n}{b} \PY{o+ow}{in} \PY{n}{combinations}\PY{p}{(}\PY{n}{G}\PY{o}{.}\PY{n}{nodes}\PY{p}{(}\PY{p}{)}\PY{p}{,} \PY{l+m+mi}{2}\PY{p}{)}\PY{p}{:}\PY{c+c1}{\PYZsh{}For all \PYZdq{}n choose 2\PYZdq{} pairs of nodes...}
        \PY{k}{if} \PY{n}{node\PYZus{}i} \PY{o+ow}{in} \PY{p}{(}\PY{n}{a}\PY{p}{,} \PY{n}{b}\PY{p}{)}\PY{p}{:} \PY{c+c1}{\PYZsh{}Excluding cases where node\PYZus{}i is one of those nodes.}
            \PY{k}{continue}
        \PY{k}{try}\PY{p}{:}
            \PY{n}{paths} \PY{o}{=} \PY{n+nb}{list}\PY{p}{(}\PY{n}{nx}\PY{o}{.}\PY{n}{all\PYZus{}shortest\PYZus{}paths}\PY{p}{(}\PY{n}{G}\PY{p}{,} \PY{n}{a}\PY{p}{,} \PY{n}{b}\PY{p}{)}\PY{p}{)} \PY{c+c1}{\PYZsh{}... compute all possible shortest paths from this pair}
        \PY{k}{except} \PY{n}{nx}\PY{o}{.}\PY{n}{NetworkXNoPath}\PY{p}{:}
            \PY{k}{continue}

        \PY{n}{g\PYZus{}ab} \PY{o}{=} \PY{n+nb}{len}\PY{p}{(}\PY{n}{paths}\PY{p}{)}\PY{c+c1}{\PYZsh{}The total number of geodesics}
        \PY{n}{g\PYZus{}ab\PYZus{}i} \PY{o}{=} \PY{l+m+mi}{0}\PY{c+c1}{\PYZsh{}The number of geodesics with node\PYZus{}i in them}

        \PY{k}{for} \PY{n}{p} \PY{o+ow}{in} \PY{n}{paths}\PY{p}{:} \PY{c+c1}{\PYZsh{}For all geodesics (shortest paths)...}
            \PY{k}{if} \PY{n}{node\PYZus{}i} \PY{o+ow}{in} \PY{n}{p}\PY{p}{:}\PY{c+c1}{\PYZsh{}...if node\PYZus{}i is in a geodesic...}
                \PY{n}{g\PYZus{}ab\PYZus{}i} \PY{o}{+}\PY{o}{=} \PY{l+m+mi}{1} \PY{c+c1}{\PYZsh{}... count it!}

        \PY{n}{centrality} \PY{o}{+}\PY{o}{=} \PY{n}{g\PYZus{}ab\PYZus{}i} \PY{o}{/} \PY{n}{g\PYZus{}ab} \PY{c+c1}{\PYZsh{}C\PYZus{}B(node\PYZus{}i), the sum, of the proportion of geodesics node\PYZus{}i is a part of, for all \PYZdq{}n choose 2\PYZdq{} combinations}

    \PY{n}{max\PYZus{}centrality} \PY{o}{=} \PY{p}{(}\PY{p}{(}\PY{n}{N}\PY{o}{*}\PY{o}{*}\PY{l+m+mi}{2}\PY{p}{)}\PY{o}{\PYZhy{}}\PY{p}{(}\PY{l+m+mi}{3}\PY{o}{*}\PY{n}{N}\PY{p}{)}\PY{o}{+}\PY{l+m+mi}{2}\PY{p}{)}\PY{o}{/}\PY{l+m+mi}{2} \PY{c+c1}{\PYZsh{} max possible betweenness centrality}
    \PY{n}{centrality} \PY{o}{=} \PY{n}{centrality}\PY{o}{/}\PY{n}{max\PYZus{}centrality} \PY{c+c1}{\PYZsh{} betweeness centrality of node\PYZus{}i relative to the max possible betweeness centrality of \PYZsq{}G\PYZsq{}}
    \PY{k}{return} \PY{n}{centrality}
\end{Verbatim}
\end{tcolorbox}

    \begin{tcolorbox}[breakable, size=fbox, boxrule=1pt, pad at break*=1mm,colback=cellbackground, colframe=cellborder]
\prompt{In}{incolor}{6}{\boxspacing}
\begin{Verbatim}[commandchars=\\\{\}]
\PY{k+kn}{import}\PY{+w}{ }\PY{n+nn}{networkx}\PY{+w}{ }\PY{k}{as}\PY{+w}{ }\PY{n+nn}{nx}
\PY{k+kn}{import}\PY{+w}{ }\PY{n+nn}{matplotlib}\PY{n+nn}{.}\PY{n+nn}{pyplot}\PY{+w}{ }\PY{k}{as}\PY{+w}{ }\PY{n+nn}{plt}
\PY{n}{kcg} \PY{o}{=} \PY{n}{nx}\PY{o}{.}\PY{n}{karate\PYZus{}club\PYZus{}graph}\PY{p}{(}\PY{p}{)}
\PY{n+nb}{print}\PY{p}{(}\PY{n}{kcg}\PY{o}{.}\PY{n+nv+vm}{\PYZus{}\PYZus{}class\PYZus{}\PYZus{}}\PY{p}{)}
\PY{n}{betweenness\PYZus{}centrality} \PY{o}{=} \PY{p}{\PYZob{}}\PY{p}{\PYZcb{}}
\PY{k}{for} \PY{n}{node\PYZus{}i} \PY{o+ow}{in} \PY{n}{kcg}\PY{o}{.}\PY{n}{nodes}\PY{p}{(}\PY{p}{)}\PY{p}{:}
    \PY{n}{betweenness\PYZus{}centrality}\PY{p}{[}\PY{n}{node\PYZus{}i}\PY{p}{]} \PY{o}{=} \PY{n}{betweenness\PYZus{}centrality\PYZus{}from\PYZus{}scratch}\PY{p}{(}\PY{n}{kcg}\PY{p}{,} \PY{n}{node\PYZus{}i}\PY{p}{)}
\PY{n}{betweenness\PYZus{}centrality} \PY{o}{=} \PY{p}{(}\PY{n+nb}{sorted}\PY{p}{(}\PY{n}{betweenness\PYZus{}centrality}\PY{o}{.}\PY{n}{items}\PY{p}{(}\PY{p}{)}\PY{p}{,} \PY{n}{key}\PY{o}{=}\PY{k}{lambda} \PY{n}{x}\PY{p}{:} \PY{n}{x}\PY{p}{[}\PY{l+m+mi}{1}\PY{p}{]}\PY{p}{,} \PY{n}{reverse}\PY{o}{=}\PY{k+kc}{True}\PY{p}{)}\PY{p}{)}
\PY{n+nb}{print}\PY{p}{(}\PY{l+s+s2}{\PYZdq{}}\PY{l+s+s2}{My Version: }\PY{l+s+s2}{\PYZdq{}}\PY{p}{)}
\PY{n+nb}{print}\PY{p}{(}\PY{n}{betweenness\PYZus{}centrality}\PY{p}{)}
\PY{n+nb}{print}\PY{p}{(}\PY{l+s+s2}{\PYZdq{}}\PY{l+s+s2}{Networkx Version: }\PY{l+s+s2}{\PYZdq{}}\PY{p}{)}
\PY{n}{betweenness\PYZus{}centrality\PYZus{}official} \PY{o}{=} \PY{p}{(}\PY{n+nb}{sorted}\PY{p}{(}\PY{n}{nx}\PY{o}{.}\PY{n}{betweenness\PYZus{}centrality}\PY{p}{(}\PY{n}{kcg}\PY{p}{,}\PY{p}{)}\PY{o}{.}\PY{n}{items}\PY{p}{(}\PY{p}{)}\PY{p}{,} \PY{n}{key}\PY{o}{=}\PY{k}{lambda} \PY{n}{x}\PY{p}{:} \PY{n}{x}\PY{p}{[}\PY{l+m+mi}{1}\PY{p}{]}\PY{p}{,} \PY{n}{reverse}\PY{o}{=}\PY{k+kc}{True}\PY{p}{)}\PY{p}{)}
\PY{n+nb}{print}\PY{p}{(}\PY{n}{betweenness\PYZus{}centrality\PYZus{}official}\PY{p}{)}

\PY{n}{pos} \PY{o}{=} \PY{n}{nx}\PY{o}{.}\PY{n}{spring\PYZus{}layout}\PY{p}{(}\PY{n}{kcg}\PY{p}{)}
\PY{n}{plt}\PY{o}{.}\PY{n}{plot}\PY{p}{(}\PY{p}{)}
\PY{n}{nx}\PY{o}{.}\PY{n}{draw}\PY{p}{(}\PY{n}{kcg}\PY{p}{,} \PY{n}{pos}\PY{p}{,} \PY{n}{node\PYZus{}size}\PY{o}{=}\PY{l+m+mi}{150}\PY{p}{,} \PY{n}{edgecolors}\PY{o}{=}\PY{l+s+s1}{\PYZsq{}}\PY{l+s+s1}{.2}\PY{l+s+s1}{\PYZsq{}}\PY{p}{,}
        \PY{n}{with\PYZus{}labels}\PY{o}{=}\PY{k+kc}{True}\PY{p}{,} \PY{n}{font\PYZus{}size}\PY{o}{=}\PY{l+s+s1}{\PYZsq{}}\PY{l+s+s1}{x\PYZhy{}small}\PY{l+s+s1}{\PYZsq{}}\PY{p}{,} \PY{n}{edge\PYZus{}color}\PY{o}{=}\PY{l+s+s1}{\PYZsq{}}\PY{l+s+s1}{.6}\PY{l+s+s1}{\PYZsq{}}\PY{p}{)}
\end{Verbatim}
\end{tcolorbox}

    \begin{Verbatim}[commandchars=\\\{\}]
<class 'networkx.classes.graph.Graph'>
My Version:
[(0, 0.4376352813852816), (33, 0.3040749759499759), (32, 0.14524711399711399),
(2, 0.1436568061568063), (31, 0.13827561327561325), (8, 0.05592682780182779),
(1, 0.0539366883116883), (13, 0.045863395863395835), (19, 0.03247504810004808),
(5, 0.02998737373737374), (6, 0.02998737373737374), (27, 0.022333453583453584),
(23, 0.01761363636363637), (30, 0.014411976911976923), (3,
0.011909271284271286), (25, 0.0038404882154882154), (29, 0.0029220779220779218),
(24, 0.0022095959595959595), (28, 0.0017947330447330445), (9,
0.0008477633477633477), (4, 0.0006313131313131313), (10, 0.0006313131313131313),
(7, 0.0), (11, 0.0), (12, 0.0), (14, 0.0), (15, 0.0), (16, 0.0), (17, 0.0), (18,
0.0), (20, 0.0), (21, 0.0), (22, 0.0), (26, 0.0)]
Networkx Version:
[(0, 0.43763528138528146), (33, 0.30407497594997596), (32, 0.145247113997114),
(2, 0.14365680615680618), (31, 0.13827561327561325), (8, 0.05592682780182781),
(1, 0.053936688311688304), (13, 0.04586339586339586), (19, 0.03247504810004811),
(5, 0.02998737373737374), (6, 0.029987373737373736), (27, 0.02233345358345358),
(23, 0.017613636363636363), (30, 0.014411976911976909), (3,
0.011909271284271283), (25, 0.0038404882154882154), (29, 0.0029220779220779218),
(24, 0.0022095959595959595), (28, 0.0017947330447330447), (9,
0.0008477633477633478), (4, 0.0006313131313131313), (10, 0.0006313131313131313),
(7, 0.0), (11, 0.0), (12, 0.0), (14, 0.0), (15, 0.0), (16, 0.0), (17, 0.0), (18,
0.0), (20, 0.0), (21, 0.0), (22, 0.0), (26, 0.0)]
    \end{Verbatim}

    \begin{center}
    \adjustimage{max size={0.9\linewidth}{0.9\paperheight}}{Hughes_assignment01_files/Hughes_assignment01_23_1.png}
    \end{center}
    { \hspace*{\fill} \\}
    
    \textbf{(b) Augment your function with error handling. It should be able
to return useful} \textbf{messages when given any of the following:} *
i. A node that is not in the graph * ii. A graph that is not a networkx
object * iii. A graph with no connectivity (or not enough connectivity
to produce a valid centrality result) * iv. Any situation that might
produce numerical difficulties (e.g.~dividing by zero) for your chosen
centrality measure

    \begin{tcolorbox}[breakable, size=fbox, boxrule=1pt, pad at break*=1mm,colback=cellbackground, colframe=cellborder]
\prompt{In}{incolor}{19}{\boxspacing}
\begin{Verbatim}[commandchars=\\\{\}]
\PY{k}{def}\PY{+w}{ }\PY{n+nf}{betweenness\PYZus{}centrality\PYZus{}from\PYZus{}scratch}\PY{p}{(}\PY{n}{G}\PY{p}{,} \PY{n}{node\PYZus{}i}\PY{p}{)}\PY{p}{:}
\PY{+w}{    }\PY{l+s+sd}{\PYZdq{}\PYZdq{}\PYZdq{}}
\PY{l+s+sd}{    Calculate the betweenness centrality for each node in a graph from scratch.}

\PY{l+s+sd}{    Betweenness centrality is a measure of the extent to which a node lies on the shortest}
\PY{l+s+sd}{    paths between other nodes in the network. This function computes the betweenness}
\PY{l+s+sd}{    centrality for all nodes in the graph `G`.}

\PY{l+s+sd}{    Parameters}
\PY{l+s+sd}{    \PYZhy{}\PYZhy{}\PYZhy{}\PYZhy{}\PYZhy{}\PYZhy{}\PYZhy{}\PYZhy{}\PYZhy{}\PYZhy{}}
\PY{l+s+sd}{    G : networkx.Graph}
\PY{l+s+sd}{        The input graph on which the betweenness centrality is calculated. It can}
\PY{l+s+sd}{        be any type of graph (undirected, directed, etc.) supported by NetworkX.}

\PY{l+s+sd}{    node\PYZus{}i : networkx.Node}
\PY{l+s+sd}{        The input node on which the betweenness centrality is calculated.}

\PY{l+s+sd}{    Returns}
\PY{l+s+sd}{    \PYZhy{}\PYZhy{}\PYZhy{}\PYZhy{}\PYZhy{}\PYZhy{}\PYZhy{}}
\PY{l+s+sd}{    centrality : float}
\PY{l+s+sd}{        A float where the value is the corresponding betweenness centrality score as a function of the total number of nodes N, the total number of geodesics for all pairs of nodes in  graph G, and the number of geodesics node\PYZus{}i is part of in said graph, G.}


\PY{l+s+sd}{    Time Complexity}
\PY{l+s+sd}{    \PYZhy{}\PYZhy{}\PYZhy{}\PYZhy{}\PYZhy{}\PYZhy{}\PYZhy{}\PYZhy{}\PYZhy{}\PYZhy{}\PYZhy{}\PYZhy{}\PYZhy{}\PYZhy{}\PYZhy{}}
\PY{l+s+sd}{    The overall time complexity is O(N\PYZca{}2) where N is the number of nodes, due to the use of}
\PY{l+s+sd}{    combinations(G.nodes(), 2) dominating, which provides all \PYZdq{}N choose 2\PYZdq{} pairs of nodes.}

\PY{l+s+sd}{    Notes}
\PY{l+s+sd}{    \PYZhy{}\PYZhy{}\PYZhy{}\PYZhy{}\PYZhy{}}
\PY{l+s+sd}{    If unweighted, networkx\PYZsq{}s all\PYZus{}shortest\PYZus{}paths() uses Breadth\PYZhy{}First Search which is O(N), if weighted though, nx.all\PYZus{}shortest\PYZus{}paths() defaults to}
\PY{l+s+sd}{    Dijkstra\PYZsq{}s which is O(E*N*log(N)), and if negative weights exist then it is O(N*E) due to Bellman\PYZhy{}Ford. All where N is the number of nodes, where E is}
\PY{l+s+sd}{    the number of edges.}
\PY{l+s+sd}{    https://networkx.org/documentation/stable/\PYZus{}modules/networkx/algorithms/shortest\PYZus{}paths/unweighted.html}

\PY{l+s+sd}{    Citations}
\PY{l+s+sd}{    \PYZhy{}\PYZhy{}\PYZhy{}\PYZhy{}\PYZhy{}\PYZhy{}\PYZhy{}\PYZhy{}\PYZhy{}}
\PY{l+s+sd}{    Freeman, Linton C. “A Set of Measures of Centrality Based on Betweenness.” Sociometry, vol. 40, no. 1, 1977, pp. 35–41. JSTOR, https://doi.org/10.2307/3033543. Accessed 6 Oct. 2026.}

\PY{l+s+sd}{    Example}
\PY{l+s+sd}{    \PYZhy{}\PYZhy{}\PYZhy{}\PYZhy{}\PYZhy{}\PYZhy{}\PYZhy{}}
\PY{l+s+sd}{    \PYZgt{}\PYZgt{}\PYZgt{} import networkx as nx}
\PY{l+s+sd}{    \PYZgt{}\PYZgt{}\PYZgt{} G = nx.karate\PYZus{}club\PYZus{}graph()}
\PY{l+s+sd}{    \PYZgt{}\PYZgt{}\PYZgt{} node\PYZus{}i = 3}
\PY{l+s+sd}{    \PYZgt{}\PYZgt{}\PYZgt{} betweenness\PYZus{}centrality\PYZus{}from\PYZus{}scratch(kcg, node\PYZus{}i)}
\PY{l+s+sd}{    (3, 0.011909271284271286)}
\PY{l+s+sd}{    \PYZdq{}\PYZdq{}\PYZdq{}}
    \PY{k+kn}{import}\PY{+w}{ }\PY{n+nn}{networkx}\PY{+w}{ }\PY{k}{as}\PY{+w}{ }\PY{n+nn}{nx}
    \PY{k+kn}{from}\PY{+w}{ }\PY{n+nn}{itertools}\PY{+w}{ }\PY{k+kn}{import} \PY{n}{combinations}
    \PY{k+kn}{import}\PY{+w}{ }\PY{n+nn}{logging}

    \PY{c+c1}{\PYZsh{}Error Handling Section}
    \PY{c+c1}{\PYZsh{}Check that G is a valid graph}
    \PY{k}{if} \PY{o+ow}{not} \PY{n+nb}{isinstance}\PY{p}{(}\PY{n}{G}\PY{p}{,}\PY{n}{nx}\PY{o}{.}\PY{n}{Graph}\PY{p}{)}\PY{p}{:}
        \PY{k}{raise}  \PY{n+ne}{ValueError}\PY{p}{(}\PY{l+s+s2}{\PYZdq{}}\PY{l+s+s2}{Graph provided is not a networkx graph!}\PY{l+s+s2}{\PYZdq{}}\PY{p}{)}
    \PY{c+c1}{\PYZsh{}Node not a part of graph}
    \PY{k}{if} \PY{n}{node\PYZus{}i} \PY{o+ow}{not} \PY{o+ow}{in} \PY{n}{G}\PY{o}{.}\PY{n}{nodes}\PY{p}{(}\PY{p}{)}\PY{p}{:}
        \PY{k}{raise}  \PY{n+ne}{ValueError}\PY{p}{(}\PY{l+s+s2}{\PYZdq{}}\PY{l+s+s2}{Node provided is not in graph!}\PY{l+s+s2}{\PYZdq{}}\PY{p}{)}
    \PY{c+c1}{\PYZsh{} Not enough edges}
    \PY{k}{if} \PY{n}{G}\PY{o}{.}\PY{n}{number\PYZus{}of\PYZus{}edges}\PY{p}{(}\PY{p}{)} \PY{o}{\PYZlt{}} \PY{l+m+mi}{1}\PY{p}{:}
        \PY{k}{raise}  \PY{n+ne}{ValueError}\PY{p}{(}\PY{l+s+s2}{\PYZdq{}}\PY{l+s+s2}{Graph provided has no connectivity!}\PY{l+s+s2}{\PYZdq{}}\PY{p}{)}

    \PY{c+c1}{\PYZsh{}Catch any errors raised before we start:}
    \PY{k}{try}\PY{p}{:}
        \PY{n}{centrality} \PY{o}{=} \PY{l+m+mf}{0.0}
        \PY{n}{N} \PY{o}{=} \PY{n}{G}\PY{o}{.}\PY{n}{number\PYZus{}of\PYZus{}nodes}\PY{p}{(}\PY{p}{)}  \PY{c+c1}{\PYZsh{} Total number of nodes in the graph}
    \PY{k}{except} \PY{n+ne}{ValueError}\PY{p}{:}
        \PY{n}{logging}\PY{o}{.}\PY{n}{debug}\PY{p}{(}\PY{l+s+s1}{\PYZsq{}}\PY{l+s+s1}{Value Error: Invalid Input}\PY{l+s+s1}{\PYZsq{}}\PY{p}{)}
    \PY{k}{except} \PY{n+ne}{Exception} \PY{k}{as} \PY{n}{e}\PY{p}{:}
            \PY{n+nb}{print}\PY{p}{(}\PY{n}{e}\PY{p}{)}
            \PY{n}{logging}\PY{o}{.}\PY{n}{debug}\PY{p}{(}\PY{l+s+s1}{\PYZsq{}}\PY{l+s+s1}{other error}\PY{l+s+s1}{\PYZsq{}}\PY{p}{)}

\PY{+w}{    }\PY{l+s+sd}{\PYZsq{}\PYZsq{}\PYZsq{}}
\PY{l+s+sd}{    Grabbing C\PYZus{}B(p\PYZus{}k) =\PYZgt{} C\PYZus{}B(node\PYZus{}i) = \PYZdq{}the overall partial betweenness of point p\PYZus{}k.\PYZdq{}}
\PY{l+s+sd}{    If there is only 1 geodesic (i.e., shortest path) C\PYZus{}B(p\PYZus{}k) is increased by 1.}
\PY{l+s+sd}{    Otherwise, C\PYZus{}B(p\PYZus{}k) is increased relative to the proportion of geodesics p\PYZus{}k is apart of}
\PY{l+s+sd}{    \PYZsq{}\PYZsq{}\PYZsq{}}
    \PY{k}{for} \PY{n}{a}\PY{p}{,} \PY{n}{b} \PY{o+ow}{in} \PY{n}{combinations}\PY{p}{(}\PY{n}{G}\PY{o}{.}\PY{n}{nodes}\PY{p}{(}\PY{p}{)}\PY{p}{,} \PY{l+m+mi}{2}\PY{p}{)}\PY{p}{:}              \PY{c+c1}{\PYZsh{}For all \PYZdq{}n choose 2\PYZdq{} pairs of nodes...}
        \PY{k}{if} \PY{n}{node\PYZus{}i} \PY{o+ow}{in} \PY{p}{(}\PY{n}{a}\PY{p}{,} \PY{n}{b}\PY{p}{)}\PY{p}{:} \PY{c+c1}{\PYZsh{}Excluding cases where node\PYZus{}i is one of those nodes.}
            \PY{k}{continue}
        \PY{k}{try}\PY{p}{:}
            \PY{n}{paths} \PY{o}{=} \PY{n+nb}{list}\PY{p}{(}\PY{n}{nx}\PY{o}{.}\PY{n}{all\PYZus{}shortest\PYZus{}paths}\PY{p}{(}\PY{n}{G}\PY{p}{,} \PY{n}{a}\PY{p}{,} \PY{n}{b}\PY{p}{)}\PY{p}{)} \PY{c+c1}{\PYZsh{}... compute all possible shortest paths from this pair}
        \PY{k}{except} \PY{n}{nx}\PY{o}{.}\PY{n}{NetworkXNoPath}\PY{p}{:}
            \PY{k}{continue}
        \PY{k}{except} \PY{n+ne}{Exception} \PY{k}{as} \PY{n}{e}\PY{p}{:}
            \PY{n+nb}{print}\PY{p}{(}\PY{n}{e}\PY{p}{)}
            \PY{n}{logging}\PY{o}{.}\PY{n}{debug}\PY{p}{(}\PY{l+s+s1}{\PYZsq{}}\PY{l+s+s1}{other error}\PY{l+s+s1}{\PYZsq{}}\PY{p}{)}

        \PY{n}{g\PYZus{}ab} \PY{o}{=} \PY{n+nb}{len}\PY{p}{(}\PY{n}{paths}\PY{p}{)}                                \PY{c+c1}{\PYZsh{}The total number of geodesics}
        \PY{n}{g\PYZus{}ab\PYZus{}i} \PY{o}{=} \PY{l+m+mi}{0}                                       \PY{c+c1}{\PYZsh{}The number of geodesics with node\PYZus{}i in them}

        \PY{k}{for} \PY{n}{p} \PY{o+ow}{in} \PY{n}{paths}\PY{p}{:} \PY{c+c1}{\PYZsh{}For all geodesics (shortest paths)...}
            \PY{k}{if} \PY{n}{node\PYZus{}i} \PY{o+ow}{in} \PY{n}{p}\PY{p}{:}\PY{c+c1}{\PYZsh{}...if node\PYZus{}i is in a geodesic...}
                \PY{n}{g\PYZus{}ab\PYZus{}i} \PY{o}{+}\PY{o}{=} \PY{l+m+mi}{1} \PY{c+c1}{\PYZsh{}... count it!}
        \PY{k}{if}\PY{p}{(}\PY{n}{g\PYZus{}ab\PYZus{}i}\PY{o}{==}\PY{l+m+mi}{0}\PY{p}{)}\PY{p}{:}
            \PY{k}{raise} \PY{n+ne}{ValueError}\PY{p}{(}\PY{l+s+s1}{\PYZsq{}}\PY{l+s+s1}{Value Error: node\PYZus{}i not on any geodesics!}\PY{l+s+s1}{\PYZsq{}}\PY{p}{)}
        \PY{k}{try}\PY{p}{:}
            \PY{n}{centrality} \PY{o}{+}\PY{o}{=} \PY{n}{g\PYZus{}ab\PYZus{}i} \PY{o}{/} \PY{n}{g\PYZus{}ab} \PY{c+c1}{\PYZsh{}C\PYZus{}B(node\PYZus{}i), the sum, of the proportion of geodesics node\PYZus{}i is a part of, for all \PYZdq{}n choose 2\PYZdq{} combinations}
        \PY{k}{except} \PY{n+ne}{ValueError}\PY{p}{:}
            \PY{n}{logging}\PY{o}{.}\PY{n}{debug}\PY{p}{(}\PY{l+s+s1}{\PYZsq{}}\PY{l+s+s1}{Value Error: Invalid Input}\PY{l+s+s1}{\PYZsq{}}\PY{p}{)}
        \PY{k}{except} \PY{n+ne}{Exception} \PY{k}{as} \PY{n}{e}\PY{p}{:}
                \PY{n+nb}{print}\PY{p}{(}\PY{n}{e}\PY{p}{)}
                \PY{n}{logging}\PY{o}{.}\PY{n}{debug}\PY{p}{(}\PY{l+s+s1}{\PYZsq{}}\PY{l+s+s1}{other error}\PY{l+s+s1}{\PYZsq{}}\PY{p}{)}

    \PY{n}{max\PYZus{}centrality} \PY{o}{=} \PY{p}{(}\PY{p}{(}\PY{n}{N}\PY{o}{*}\PY{o}{*}\PY{l+m+mi}{2}\PY{p}{)}\PY{o}{\PYZhy{}}\PY{p}{(}\PY{l+m+mi}{3}\PY{o}{*}\PY{n}{N}\PY{p}{)}\PY{o}{+}\PY{l+m+mi}{2}\PY{p}{)}\PY{o}{/}\PY{l+m+mi}{2} \PY{c+c1}{\PYZsh{} max possible betweenness centrality}
    \PY{n}{centrality} \PY{o}{=} \PY{n}{centrality}\PY{o}{/}\PY{n}{max\PYZus{}centrality} \PY{c+c1}{\PYZsh{} betweeness centrality of node\PYZus{}i relative to the max possible betweeness centrality of \PYZsq{}G\PYZsq{}}
    \PY{k}{return} \PY{n}{centrality}
\end{Verbatim}
\end{tcolorbox}

    \begin{enumerate}
\def\labelenumi{(\alph{enumi})}
\setcounter{enumi}{2}
\tightlist
\item
  Compare your function to the corresponding networkx function across
  several graphs you design as test inputs, including one with no
  connectivity. Does your function ever return different results from
  the networkx function? Why do you think this is the case?
\end{enumerate}

    \emph{In my first attempt, my version only aligned relatively speaking,
where I produced slightly different values and orders than then one in
networkX. E.g., I calculated node 2 to be the 2nd most central while
networkx did so for node 33. Now however, I have it fixed and I believe
it aligns fully!}

\emph{Originally I didn't think of cases where there may be multiple
equally valid shortest paths between 2 nodes, where node\_i is only on a
portion of said valid list of shortest paths.}

    \begin{enumerate}
\def\labelenumi{(\alph{enumi})}
\setcounter{enumi}{3}
\tightlist
\item
  Putting it all together: Write a function that takes a centrality
  function as input and tests the function for robustness to edge cases
  (like a graph with no connectivity) and correctness of outputs.
  Explain why you selected the test cases you chose.
\end{enumerate}

    \begin{tcolorbox}[breakable, size=fbox, boxrule=1pt, pad at break*=1mm,colback=cellbackground, colframe=cellborder]
\prompt{In}{incolor}{8}{\boxspacing}
\begin{Verbatim}[commandchars=\\\{\}]
\PY{k+kn}{import}\PY{+w}{ }\PY{n+nn}{networkx}\PY{+w}{ }\PY{k}{as}\PY{+w}{ }\PY{n+nn}{nx}
\PY{k+kn}{import}\PY{+w}{ }\PY{n+nn}{numpy}\PY{+w}{ }\PY{k}{as}\PY{+w}{ }\PY{n+nn}{np}
\PY{n}{kcg} \PY{o}{=} \PY{n}{nx}\PY{o}{.}\PY{n}{karate\PYZus{}club\PYZus{}graph}\PY{p}{(}\PY{p}{)}
\PY{n}{node\PYZus{}i} \PY{o}{=} \PY{l+m+mi}{3}
\PY{n+nb}{print}\PY{p}{(}\PY{n}{kcg}\PY{o}{.}\PY{n+nv+vm}{\PYZus{}\PYZus{}class\PYZus{}\PYZus{}}\PY{p}{)}
\PY{n}{betweenness\PYZus{}centrality} \PY{o}{=} \PY{p}{\PYZob{}}\PY{p}{\PYZcb{}}


\PY{c+c1}{\PYZsh{}Check valid graph}
\PY{n}{betweenness\PYZus{}centrality\PYZus{}from\PYZus{}scratch}\PY{p}{(}\PY{l+m+mi}{1}\PY{p}{,} \PY{n}{node\PYZus{}i}\PY{p}{)}
\PY{c+c1}{\PYZsh{}Check valid node}
\PY{n}{betweenness\PYZus{}centrality\PYZus{}from\PYZus{}scratch}\PY{p}{(}\PY{n}{kcg}\PY{p}{,} \PY{l+m+mi}{50}\PY{p}{)}


\PY{n}{G} \PY{o}{=} \PY{n}{nx}\PY{o}{.}\PY{n}{Graph}\PY{p}{(}\PY{p}{)}
\PY{n}{G}\PY{o}{.}\PY{n}{add\PYZus{}nodes\PYZus{}from}\PY{p}{(}\PY{n+nb}{reversed}\PY{p}{(}\PY{n+nb}{range}\PY{p}{(}\PY{l+m+mi}{10}\PY{p}{)}\PY{p}{)}\PY{p}{)}
\PY{n+nb}{print}\PY{p}{(}\PY{n+nb}{list}\PY{p}{(}\PY{n}{G}\PY{o}{.}\PY{n}{subgraph}\PY{p}{(}\PY{p}{[}\PY{l+m+mi}{1}\PY{p}{,}\PY{l+m+mi}{2}\PY{p}{,}\PY{l+m+mi}{3}\PY{p}{]}\PY{p}{)}\PY{p}{)}\PY{p}{)}

\PY{c+c1}{\PYZsh{}Check valid connectivity}
\PY{n}{node\PYZus{}i}\PY{o}{=}\PY{l+m+mi}{1}
\PY{n}{betweenness\PYZus{}centrality\PYZus{}from\PYZus{}scratch}\PY{p}{(}\PY{n}{G}\PY{o}{.}\PY{n}{subgraph}\PY{p}{(}\PY{p}{[}\PY{l+m+mi}{1}\PY{p}{]}\PY{p}{)}\PY{p}{,} \PY{n}{node\PYZus{}i}\PY{p}{)}
\end{Verbatim}
\end{tcolorbox}

    \begin{Verbatim}[commandchars=\\\{\}]
<class 'networkx.classes.graph.Graph'>
    \end{Verbatim}

    \begin{Verbatim}[commandchars=\\\{\}, frame=single, framerule=2mm, rulecolor=\color{outerrorbackground}]
\textcolor{ansi-red}{---------------------------------------------------------------------------}
\textcolor{ansi-red}{ValueError}                                Traceback (most recent call last)
\textcolor{ansi-cyan}{Cell}\textcolor{ansi-cyan}{ }\textcolor{ansi-green}{In[8]}\textcolor{ansi-green}{, line 10}
\textcolor{ansi-green}{      6} betweenness\_centrality = \{\}
\textcolor{ansi-green}{      7} 
\textcolor{ansi-green}{      8} 
\textcolor{ansi-green}{      9} \def\tcRGB{\textcolor[RGB]}\expandafter\tcRGB\expandafter{\detokenize{95,135,135}}{\#Check valid graph}
\textcolor{ansi-green}{---> }\textcolor{ansi-green}{10} betweenness\_centrality\_from\_scratch(\textcolor{ansi-green}{1}, node\_i)
\textcolor{ansi-green}{     11} \def\tcRGB{\textcolor[RGB]}\expandafter\tcRGB\expandafter{\detokenize{95,135,135}}{\#Check valid node}
\textcolor{ansi-green}{     12} betweenness\_centrality\_from\_scratch(kcg, \textcolor{ansi-green}{50})
\textcolor{ansi-green}{     13} 

\textcolor{ansi-cyan}{Cell}\textcolor{ansi-cyan}{ }\textcolor{ansi-green}{In[7]}\textcolor{ansi-green}{, line 55}, in \textcolor{ansi-cyan}{betweenness\_centrality\_from\_scratch}\textcolor{ansi-blue}{(G, node\_i)}
\textcolor{ansi-green}{     51} 
\textcolor{ansi-green}{     52}     \def\tcRGB{\textcolor[RGB]}\expandafter\tcRGB\expandafter{\detokenize{95,135,135}}{\#Error Handling Section}
\textcolor{ansi-green}{     53}     \def\tcRGB{\textcolor[RGB]}\expandafter\tcRGB\expandafter{\detokenize{95,135,135}}{\#Check that G is a valid graph}
\textcolor{ansi-green}{     54}     \def\tcRGB{\textcolor[RGB]}\expandafter\tcRGB\expandafter{\detokenize{0,135,0}}{\textbf{if}} \def\tcRGB{\textcolor[RGB]}\expandafter\tcRGB\expandafter{\detokenize{0,135,0}}{\textbf{not}} isinstance(G,nx.Graph):
\textcolor{ansi-green}{---> }\textcolor{ansi-green}{55}         \def\tcRGB{\textcolor[RGB]}\expandafter\tcRGB\expandafter{\detokenize{0,135,0}}{\textbf{raise}}  ValueError(\textcolor{ansi-yellow}{"Graph provided is not a networkx graph!"})
\textcolor{ansi-green}{     56}     \def\tcRGB{\textcolor[RGB]}\expandafter\tcRGB\expandafter{\detokenize{95,135,135}}{\#Node not a part of graph}
\textcolor{ansi-green}{     57}     \def\tcRGB{\textcolor[RGB]}\expandafter\tcRGB\expandafter{\detokenize{0,135,0}}{\textbf{if}} node\_i \def\tcRGB{\textcolor[RGB]}\expandafter\tcRGB\expandafter{\detokenize{0,135,0}}{\textbf{not}} \def\tcRGB{\textcolor[RGB]}\expandafter\tcRGB\expandafter{\detokenize{0,135,0}}{\textbf{in}} G.nodes():
\textcolor{ansi-green}{     58}         \def\tcRGB{\textcolor[RGB]}\expandafter\tcRGB\expandafter{\detokenize{0,135,0}}{\textbf{raise}}  ValueError(\textcolor{ansi-yellow}{"Node provided is not in graph!"})

\textcolor{ansi-red}{ValueError}: Graph provided is not a networkx graph!
    \end{Verbatim}

    I picked 3 kinds of test cases and therefore errors to handle with my
code. I focused in on input based errors: namely the Graph not being a
valid nx.Graph() object, the Graph being valid but enough data to
perform the centrality measure, and the node of interest not being a
part of the graph, or alternatively not being within any of the shortest
paths. I did these 3 cases as I figured this is where the bulk of the
edge cases may lie, and as a back up, in case an exception I did not
think of shows up, I attempt to catch for exceptions writ-large. This
helps give more aid to the user of my function such that it is easier to
use and debug code that leverages it. It also ensures a level of
robustness and reliability via the back-up general exception.

    \paragraph{Question 4 (20 points).}\label{question-4-20-points.}

Download a network dataset from https://networkrepository.com/ with
between 1, 000 \textless{} N \textless{} 20, 000 nodes, and explore it
as described below.

I am using the following dataset on Facebook Artist Pages: Rossi, R., \&
Ahmed, N. (2015). The Network Data Repository with Interactive Graph
Analytics and Visualization. Proceedings of the AAAI Conference on
Artificial Intelligence, 29(1). https://doi.org/10.1609/aaai.v29i1.9277

    \begin{tcolorbox}[breakable, size=fbox, boxrule=1pt, pad at break*=1mm,colback=cellbackground, colframe=cellborder]
\prompt{In}{incolor}{9}{\boxspacing}
\begin{Verbatim}[commandchars=\\\{\}]
\PY{k+kn}{import}\PY{+w}{ }\PY{n+nn}{networkx}\PY{+w}{ }\PY{k}{as}\PY{+w}{ }\PY{n+nn}{nx}

\PY{n}{friendships} \PY{o}{=}\PY{n}{nx}\PY{o}{.}\PY{n}{read\PYZus{}edgelist}\PY{p}{(}\PY{l+s+s2}{\PYZdq{}}\PY{l+s+s2}{data/soc\PYZhy{}hamsterster/soc\PYZhy{}hamsterster.edges}\PY{l+s+s2}{\PYZdq{}}\PY{p}{,}\PY{n}{delimiter}\PY{o}{=}\PY{l+s+s2}{\PYZdq{}}\PY{l+s+s2}{ }\PY{l+s+s2}{\PYZdq{}}\PY{p}{)}
\end{Verbatim}
\end{tcolorbox}

    \begin{enumerate}
\def\labelenumi{(\alph{enumi})}
\tightlist
\item
  In 2-5 sentences, describe your network. What are the nodes, what are
  the links, what domain is it from? Be sure to include appropriate
  citation(s) to the original dataset or research paper, if relevant.
\end{enumerate}

    My network consists of friendships from the hamsterster website (Rossi,
R., \& Ahmed, N. (2015). The Network Data Repository with Interactive
Graph Analytics and Visualization. Proceedings of the AAAI Conference on
Artificial Intelligence, 29(1).
https://doi.org/10.1609/aaai.v29i1.9277). According to the documentation
found at https://networkrepository.com/soc-hamsterster.php, the network
reflects friendships and family links between users on the website.
Where each node is a user, and each edge a frienship.

    \begin{enumerate}
\def\labelenumi{(\alph{enumi})}
\setcounter{enumi}{1}
\tightlist
\item
  Compare your network's properties to those of a randomized network, by
  selecting a network null model to compare against. Your null model
  should preserve the degree distribution of the network, so randomize
  your network using e.g.~a configuration model or degree-preserving
  network rewiring (see Class 4).
\end{enumerate}

\begin{itemize}
\tightlist
\item
  \begin{enumerate}
  \def\labelenumi{\roman{enumi}.}
  \tightlist
  \item
    Create a figure with three subplots, corresponding to three
    distributions of your network; write brief interpretations of what
    you observe. These can be distributions of node properties
    (i.e.~degree, centrality measures, etc.) or edge properties
    (i.e.~edge weights). Be sure to bin these distributions
    appropriately.
  \end{enumerate}
\item
  ii. Create an analogous figure for the randomized version of your
  network; write brief interpretations of what you observe.
\end{itemize}

    \begin{tcolorbox}[breakable, size=fbox, boxrule=1pt, pad at break*=1mm,colback=cellbackground, colframe=cellborder]
\prompt{In}{incolor}{11}{\boxspacing}
\begin{Verbatim}[commandchars=\\\{\}]
\PY{k+kn}{import}\PY{+w}{ }\PY{n+nn}{networkx}\PY{+w}{ }\PY{k}{as}\PY{+w}{ }\PY{n+nn}{nx}
\PY{k+kn}{import}\PY{+w}{ }\PY{n+nn}{numpy}\PY{+w}{ }\PY{k}{as}\PY{+w}{ }\PY{n+nn}{np}
\PY{n}{N} \PY{o}{=} \PY{n}{friendships}\PY{o}{.}\PY{n}{number\PYZus{}of\PYZus{}nodes}\PY{p}{(}\PY{p}{)}
\PY{n}{L} \PY{o}{=} \PY{n}{friendships}\PY{o}{.}\PY{n}{number\PYZus{}of\PYZus{}edges}\PY{p}{(}\PY{p}{)}
\PY{n+nb}{print}\PY{p}{(}\PY{l+s+s2}{\PYZdq{}}\PY{l+s+s2}{N: }\PY{l+s+s2}{\PYZdq{}}\PY{p}{,} \PY{n}{N}\PY{p}{)}
\PY{n+nb}{print}\PY{p}{(}\PY{l+s+s2}{\PYZdq{}}\PY{l+s+s2}{L: }\PY{l+s+s2}{\PYZdq{}}\PY{p}{,} \PY{n}{L}\PY{p}{)}
\end{Verbatim}
\end{tcolorbox}

    \begin{Verbatim}[commandchars=\\\{\}]
N:  2426
L:  16630
    \end{Verbatim}

    \begin{tcolorbox}[breakable, size=fbox, boxrule=1pt, pad at break*=1mm,colback=cellbackground, colframe=cellborder]
\prompt{In}{incolor}{12}{\boxspacing}
\begin{Verbatim}[commandchars=\\\{\}]
\PY{n}{deg} \PY{o}{=} \PY{n+nb}{dict}\PY{p}{(}\PY{n}{friendships}\PY{o}{.}\PY{n}{degree}\PY{p}{(}\PY{p}{)}\PY{p}{)}
\PY{n}{avg\PYZus{}deg} \PY{o}{=} \PY{n}{np}\PY{o}{.}\PY{n}{mean}\PY{p}{(}\PY{n+nb}{list}\PY{p}{(}\PY{n}{deg}\PY{o}{.}\PY{n}{values}\PY{p}{(}\PY{p}{)}\PY{p}{)}\PY{p}{)}
\PY{n+nb}{print}\PY{p}{(}\PY{l+s+s2}{\PYZdq{}}\PY{l+s+s2}{avg\PYZus{}deg: }\PY{l+s+s2}{\PYZdq{}}\PY{p}{,} \PY{n}{avg\PYZus{}deg}\PY{p}{)}
\end{Verbatim}
\end{tcolorbox}

    \begin{Verbatim}[commandchars=\\\{\}]
avg\_deg:  13.709810387469085
    \end{Verbatim}

    \begin{tcolorbox}[breakable, size=fbox, boxrule=1pt, pad at break*=1mm,colback=cellbackground, colframe=cellborder]
\prompt{In}{incolor}{13}{\boxspacing}
\begin{Verbatim}[commandchars=\\\{\}]
\PY{n}{closeness\PYZus{}centrality} \PY{o}{=} \PY{n}{nx}\PY{o}{.}\PY{n}{closeness\PYZus{}centrality}\PY{p}{(}\PY{n}{friendships}\PY{p}{)}
\end{Verbatim}
\end{tcolorbox}

    \begin{tcolorbox}[breakable, size=fbox, boxrule=1pt, pad at break*=1mm,colback=cellbackground, colframe=cellborder]
\prompt{In}{incolor}{14}{\boxspacing}
\begin{Verbatim}[commandchars=\\\{\}]
\PY{n}{C\PYZus{}i} \PY{o}{=} \PY{n}{nx}\PY{o}{.}\PY{n}{clustering}\PY{p}{(}\PY{n}{friendships}\PY{p}{)}
\PY{n}{avg\PYZus{}C\PYZus{}i} \PY{o}{=} \PY{n}{np}\PY{o}{.}\PY{n}{mean}\PY{p}{(}\PY{n+nb}{list}\PY{p}{(}\PY{n}{C\PYZus{}i}\PY{o}{.}\PY{n}{values}\PY{p}{(}\PY{p}{)}\PY{p}{)}\PY{p}{)}
\PY{n+nb}{print}\PY{p}{(}\PY{l+s+s2}{\PYZdq{}}\PY{l+s+s2}{Average C\PYZus{}i: }\PY{l+s+s2}{\PYZdq{}}\PY{p}{,} \PY{n}{avg\PYZus{}C\PYZus{}i}\PY{p}{)}
\end{Verbatim}
\end{tcolorbox}

    \begin{Verbatim}[commandchars=\\\{\}]
Average C\_i:  0.5375333362074076
    \end{Verbatim}

    \begin{tcolorbox}[breakable, size=fbox, boxrule=1pt, pad at break*=1mm,colback=cellbackground, colframe=cellborder]
\prompt{In}{incolor}{15}{\boxspacing}
\begin{Verbatim}[commandchars=\\\{\}]
\PY{k}{def}\PY{+w}{ }\PY{n+nf}{configuration\PYZus{}model\PYZus{}from\PYZus{}degree\PYZus{}sequence}\PY{p}{(}\PY{n}{degree\PYZus{}sequence}\PY{p}{,} \PY{n}{return\PYZus{}simple}\PY{o}{=}\PY{k+kc}{True}\PY{p}{)}\PY{p}{:}
\PY{+w}{    }\PY{l+s+sd}{\PYZdq{}\PYZdq{}\PYZdq{}}
\PY{l+s+sd}{    Generate a random graph using the configuration model from a given degree sequence}
\PY{l+s+sd}{    without using the NetworkX built\PYZhy{}in function.}

\PY{l+s+sd}{    The configuration model generates a random graph that preserves the degree}
\PY{l+s+sd}{    sequence of nodes by assigning \PYZdq{}stubs\PYZdq{} or \PYZdq{}half\PYZhy{}edges\PYZdq{} to each node and}
\PY{l+s+sd}{    randomly pairing these stubs to form edges. This process can result in}
\PY{l+s+sd}{    graphs with self\PYZhy{}loops and parallel edges, which can be removed if needed.}

\PY{l+s+sd}{    Parameters}
\PY{l+s+sd}{    \PYZhy{}\PYZhy{}\PYZhy{}\PYZhy{}\PYZhy{}\PYZhy{}\PYZhy{}\PYZhy{}\PYZhy{}\PYZhy{}}
\PY{l+s+sd}{    degree\PYZus{}sequence : list of int}
\PY{l+s+sd}{        A list representing the degree of each node in the graph. The sum of}
\PY{l+s+sd}{        the degrees in this sequence must be even for the configuration model}
\PY{l+s+sd}{        to create a valid graph.}

\PY{l+s+sd}{    Returns}
\PY{l+s+sd}{    \PYZhy{}\PYZhy{}\PYZhy{}\PYZhy{}\PYZhy{}\PYZhy{}\PYZhy{}}
\PY{l+s+sd}{    G : networkx.MultiGraph}
\PY{l+s+sd}{        A multigraph generated from the given degree sequence. The graph may}
\PY{l+s+sd}{        contain self\PYZhy{}loops and parallel edges, which are allowed in the}
\PY{l+s+sd}{        configuration model.}

\PY{l+s+sd}{    Notes}
\PY{l+s+sd}{    \PYZhy{}\PYZhy{}\PYZhy{}\PYZhy{}\PYZhy{}}
\PY{l+s+sd}{    \PYZhy{} This method works by assigning \PYZdq{}stubs\PYZdq{} or \PYZdq{}half\PYZhy{}edges\PYZdq{} to nodes based on}
\PY{l+s+sd}{      their degree and then randomly pairing them to form edges. The resulting}
\PY{l+s+sd}{      graph can have self\PYZhy{}loops and parallel edges.}
\PY{l+s+sd}{    \PYZhy{} Self\PYZhy{}loops and parallel edges can be removed post\PYZhy{}generation if a simple}
\PY{l+s+sd}{      graph is required using NetworkX\PYZsq{}s `nx.Graph(G)`.}
\PY{l+s+sd}{    \PYZhy{} The degree sequence must have an even sum for a valid graph construction.}
\PY{l+s+sd}{      If the sum of the degrees is odd, no graph can be constructed.}

\PY{l+s+sd}{    Time Complexity}
\PY{l+s+sd}{    \PYZhy{}\PYZhy{}\PYZhy{}\PYZhy{}\PYZhy{}\PYZhy{}\PYZhy{}\PYZhy{}\PYZhy{}\PYZhy{}\PYZhy{}\PYZhy{}\PYZhy{}\PYZhy{}\PYZhy{}}
\PY{l+s+sd}{    The time complexity is O(E), where E is the number of edges in the graph.}

\PY{l+s+sd}{    Example}
\PY{l+s+sd}{    \PYZhy{}\PYZhy{}\PYZhy{}\PYZhy{}\PYZhy{}\PYZhy{}\PYZhy{}}
\PY{l+s+sd}{    \PYZgt{}\PYZgt{}\PYZgt{} degree\PYZus{}sequence = [3, 3, 2, 2, 1, 1]}
\PY{l+s+sd}{    \PYZgt{}\PYZgt{}\PYZgt{} G = configuration\PYZus{}model\PYZus{}from\PYZus{}degree\PYZus{}sequence(degree\PYZus{}sequence)}
\PY{l+s+sd}{    \PYZgt{}\PYZgt{}\PYZgt{} nx.is\PYZus{}graphical(degree\PYZus{}sequence)}
\PY{l+s+sd}{    True}
\PY{l+s+sd}{    \PYZdq{}\PYZdq{}\PYZdq{}}

    \PY{c+c1}{\PYZsh{} Check if the degree sequence is valid (sum of degrees must be even)}
    \PY{k}{if} \PY{n+nb}{sum}\PY{p}{(}\PY{n}{degree\PYZus{}sequence}\PY{p}{)} \PY{o}{\PYZpc{}} \PY{l+m+mi}{2} \PY{o}{!=} \PY{l+m+mi}{0}\PY{p}{:}
        \PY{k}{raise} \PY{n+ne}{ValueError}\PY{p}{(}\PY{l+s+s2}{\PYZdq{}}\PY{l+s+s2}{The sum of the degree sequence must be even.}\PY{l+s+s2}{\PYZdq{}}\PY{p}{)}

    \PY{c+c1}{\PYZsh{} Create stubs list: node i appears degree\PYZus{}sequence[i] times}
    \PY{n}{stubs} \PY{o}{=} \PY{p}{[}\PY{p}{]}
    \PY{k}{for} \PY{n}{node}\PY{p}{,} \PY{n}{degree} \PY{o+ow}{in} \PY{n+nb}{enumerate}\PY{p}{(}\PY{n}{degree\PYZus{}sequence}\PY{p}{)}\PY{p}{:}
        \PY{n}{stubs}\PY{o}{.}\PY{n}{extend}\PY{p}{(}\PY{p}{[}\PY{n}{node}\PY{p}{]} \PY{o}{*} \PY{n}{degree}\PY{p}{)}

    \PY{c+c1}{\PYZsh{} Shuffle stubs to randomize the pairing process}
    \PY{n}{np}\PY{o}{.}\PY{n}{random}\PY{o}{.}\PY{n}{shuffle}\PY{p}{(}\PY{n}{stubs}\PY{p}{)}

    \PY{c+c1}{\PYZsh{} Initialize an empty multigraph}
    \PY{n}{G} \PY{o}{=} \PY{n}{nx}\PY{o}{.}\PY{n}{MultiGraph}\PY{p}{(}\PY{p}{)}

    \PY{c+c1}{\PYZsh{} Add nodes to the graph}
    \PY{n}{G}\PY{o}{.}\PY{n}{add\PYZus{}nodes\PYZus{}from}\PY{p}{(}\PY{n+nb}{range}\PY{p}{(}\PY{n+nb}{len}\PY{p}{(}\PY{n}{degree\PYZus{}sequence}\PY{p}{)}\PY{p}{)}\PY{p}{)}

    \PY{c+c1}{\PYZsh{} Pair stubs to create edges}
    \PY{k}{while} \PY{n}{stubs}\PY{p}{:}
        \PY{n}{u} \PY{o}{=} \PY{n}{stubs}\PY{o}{.}\PY{n}{pop}\PY{p}{(}\PY{p}{)}
        \PY{n}{v} \PY{o}{=} \PY{n}{stubs}\PY{o}{.}\PY{n}{pop}\PY{p}{(}\PY{p}{)}

        \PY{c+c1}{\PYZsh{} Add the edge to the graph}
        \PY{n}{G}\PY{o}{.}\PY{n}{add\PYZus{}edge}\PY{p}{(}\PY{n}{u}\PY{p}{,} \PY{n}{v}\PY{p}{)}

    \PY{k}{if} \PY{n}{return\PYZus{}simple}\PY{p}{:}
        \PY{c+c1}{\PYZsh{} Convert the multigraph to a simple graph (remove parallel edges and self\PYZhy{}loops)}
        \PY{n}{G\PYZus{}simple} \PY{o}{=} \PY{n}{nx}\PY{o}{.}\PY{n}{Graph}\PY{p}{(}\PY{n}{G}\PY{p}{)}  \PY{c+c1}{\PYZsh{} This removes parallel edges and self\PYZhy{}loops by default}

        \PY{k}{return} \PY{n}{G\PYZus{}simple}

    \PY{k}{else}\PY{p}{:}
        \PY{k}{return} \PY{n}{G}
\end{Verbatim}
\end{tcolorbox}

    \begin{tcolorbox}[breakable, size=fbox, boxrule=1pt, pad at break*=1mm,colback=cellbackground, colframe=cellborder]
\prompt{In}{incolor}{16}{\boxspacing}
\begin{Verbatim}[commandchars=\\\{\}]
\PY{k+kn}{import}\PY{+w}{ }\PY{n+nn}{matplotlib}\PY{n+nn}{.}\PY{n+nn}{pyplot}\PY{+w}{ }\PY{k}{as}\PY{+w}{ }\PY{n+nn}{plt}

\PY{n}{fig}\PY{p}{,} \PY{n}{ax} \PY{o}{=} \PY{n}{plt}\PY{o}{.}\PY{n}{subplots}\PY{p}{(}\PY{l+m+mi}{1}\PY{p}{,}\PY{l+m+mi}{3}\PY{p}{,}\PY{n}{figsize}\PY{o}{=}\PY{p}{(}\PY{l+m+mi}{8}\PY{p}{,}\PY{l+m+mf}{3.5}\PY{p}{)}\PY{p}{,}\PY{n}{dpi}\PY{o}{=}\PY{l+m+mi}{200}\PY{p}{)}

\PY{n}{deg\PYZus{}bins} \PY{o}{=} \PY{n}{np}\PY{o}{.}\PY{n}{geomspace}\PY{p}{(}\PY{n}{np}\PY{o}{.}\PY{n}{min}\PY{p}{(}\PY{n+nb}{list}\PY{p}{(}\PY{n}{deg}\PY{o}{.}\PY{n}{values}\PY{p}{(}\PY{p}{)}\PY{p}{)}\PY{p}{)}\PY{p}{,}\PY{n}{np}\PY{o}{.}\PY{n}{max}\PY{p}{(}\PY{n+nb}{list}\PY{p}{(}\PY{n}{deg}\PY{o}{.}\PY{n}{values}\PY{p}{(}\PY{p}{)}\PY{p}{)}\PY{p}{)}\PY{p}{)}
\PY{n}{ax}\PY{p}{[}\PY{l+m+mi}{0}\PY{p}{]}\PY{o}{.}\PY{n}{hist}\PY{p}{(}\PY{n+nb}{list}\PY{p}{(}\PY{n}{deg}\PY{o}{.}\PY{n}{values}\PY{p}{(}\PY{p}{)}\PY{p}{)}\PY{p}{,}\PY{n}{bins}\PY{o}{=}\PY{n}{deg\PYZus{}bins}\PY{p}{,} \PY{n}{density}\PY{o}{=}\PY{k+kc}{True}\PY{p}{,}\PY{n}{color} \PY{o}{=} \PY{l+s+s1}{\PYZsq{}}\PY{l+s+s1}{b}\PY{l+s+s1}{\PYZsq{}}\PY{p}{,}\PY{n}{histtype}\PY{o}{=}\PY{l+s+s1}{\PYZsq{}}\PY{l+s+s1}{step}\PY{l+s+s1}{\PYZsq{}}\PY{p}{)}
\PY{n}{ax}\PY{p}{[}\PY{l+m+mi}{0}\PY{p}{]}\PY{o}{.}\PY{n}{loglog}\PY{p}{(}\PY{p}{)}
\PY{n}{ax}\PY{p}{[}\PY{l+m+mi}{0}\PY{p}{]}\PY{o}{.}\PY{n}{set\PYZus{}title}\PY{p}{(}\PY{l+s+s2}{\PYZdq{}}\PY{l+s+s2}{Degree Distribution}\PY{l+s+s2}{\PYZdq{}}\PY{p}{)}
\PY{n}{ax}\PY{p}{[}\PY{l+m+mi}{0}\PY{p}{]}\PY{o}{.}\PY{n}{set\PYZus{}xlabel}\PY{p}{(}\PY{l+s+s2}{\PYZdq{}}\PY{l+s+s2}{Degree}\PY{l+s+s2}{\PYZdq{}}\PY{p}{)}
\PY{n}{ax}\PY{p}{[}\PY{l+m+mi}{0}\PY{p}{]}\PY{o}{.}\PY{n}{set\PYZus{}ylabel}\PY{p}{(}\PY{l+s+s2}{\PYZdq{}}\PY{l+s+s2}{Probability}\PY{l+s+s2}{\PYZdq{}}\PY{p}{)}
\PY{n}{ax}\PY{p}{[}\PY{l+m+mi}{0}\PY{p}{]}\PY{o}{.}\PY{n}{grid}\PY{p}{(}\PY{k+kc}{True}\PY{p}{)}

\PY{n}{cen\PYZus{}bins} \PY{o}{=} \PY{n}{np}\PY{o}{.}\PY{n}{geomspace}\PY{p}{(}\PY{n}{np}\PY{o}{.}\PY{n}{min}\PY{p}{(}\PY{n+nb}{list}\PY{p}{(}\PY{n}{closeness\PYZus{}centrality}\PY{o}{.}\PY{n}{values}\PY{p}{(}\PY{p}{)}\PY{p}{)}\PY{p}{)}\PY{p}{,}\PY{n}{np}\PY{o}{.}\PY{n}{max}\PY{p}{(}\PY{n+nb}{list}\PY{p}{(}\PY{n}{closeness\PYZus{}centrality}\PY{o}{.}\PY{n}{values}\PY{p}{(}\PY{p}{)}\PY{p}{)}\PY{p}{)}\PY{p}{)}
\PY{n}{ax}\PY{p}{[}\PY{l+m+mi}{1}\PY{p}{]}\PY{o}{.}\PY{n}{hist}\PY{p}{(}\PY{n+nb}{list}\PY{p}{(}\PY{n}{closeness\PYZus{}centrality}\PY{o}{.}\PY{n}{values}\PY{p}{(}\PY{p}{)}\PY{p}{)}\PY{p}{,}\PY{n}{bins}\PY{o}{=}\PY{n}{cen\PYZus{}bins}\PY{p}{,}\PY{n}{density}\PY{o}{=}\PY{k+kc}{True}\PY{p}{,}\PY{n}{color} \PY{o}{=} \PY{l+s+s1}{\PYZsq{}}\PY{l+s+s1}{b}\PY{l+s+s1}{\PYZsq{}}\PY{p}{,}\PY{n}{histtype}\PY{o}{=}\PY{l+s+s1}{\PYZsq{}}\PY{l+s+s1}{step}\PY{l+s+s1}{\PYZsq{}}\PY{p}{)}
\PY{n}{ax}\PY{p}{[}\PY{l+m+mi}{1}\PY{p}{]}\PY{o}{.}\PY{n}{loglog}\PY{p}{(}\PY{p}{)}
\PY{n}{ax}\PY{p}{[}\PY{l+m+mi}{1}\PY{p}{]}\PY{o}{.}\PY{n}{set\PYZus{}title}\PY{p}{(}\PY{l+s+s2}{\PYZdq{}}\PY{l+s+s2}{Closeness Centrality Distribution}\PY{l+s+s2}{\PYZdq{}}\PY{p}{)}
\PY{n}{ax}\PY{p}{[}\PY{l+m+mi}{1}\PY{p}{]}\PY{o}{.}\PY{n}{set\PYZus{}xlabel}\PY{p}{(}\PY{l+s+s2}{\PYZdq{}}\PY{l+s+s2}{Closeness Centrality}\PY{l+s+s2}{\PYZdq{}}\PY{p}{)}
\PY{n}{ax}\PY{p}{[}\PY{l+m+mi}{1}\PY{p}{]}\PY{o}{.}\PY{n}{set\PYZus{}ylabel}\PY{p}{(}\PY{l+s+s2}{\PYZdq{}}\PY{l+s+s2}{Probability}\PY{l+s+s2}{\PYZdq{}}\PY{p}{)}
\PY{n}{ax}\PY{p}{[}\PY{l+m+mi}{1}\PY{p}{]}\PY{o}{.}\PY{n}{xaxis}\PY{o}{.}\PY{n}{set\PYZus{}inverted}\PY{p}{(}\PY{k+kc}{True}\PY{p}{)}
\PY{n}{ax}\PY{p}{[}\PY{l+m+mi}{1}\PY{p}{]}\PY{o}{.}\PY{n}{grid}\PY{p}{(}\PY{k+kc}{True}\PY{p}{)}

\PY{n}{clu\PYZus{}bins} \PY{o}{=} \PY{n}{np}\PY{o}{.}\PY{n}{geomspace}\PY{p}{(}\PY{n}{np}\PY{o}{.}\PY{n}{min}\PY{p}{(}\PY{n+nb}{list}\PY{p}{(}\PY{n}{C\PYZus{}i}\PY{o}{.}\PY{n}{values}\PY{p}{(}\PY{p}{)}\PY{p}{)}\PY{p}{)}\PY{o}{+}\PY{l+m+mf}{1e\PYZhy{}4}\PY{p}{,}\PY{n}{np}\PY{o}{.}\PY{n}{max}\PY{p}{(}\PY{n+nb}{list}\PY{p}{(}\PY{n}{C\PYZus{}i}\PY{o}{.}\PY{n}{values}\PY{p}{(}\PY{p}{)}\PY{p}{)}\PY{p}{)}\PY{p}{)}
\PY{n}{ax}\PY{p}{[}\PY{l+m+mi}{2}\PY{p}{]}\PY{o}{.}\PY{n}{hist}\PY{p}{(}\PY{n+nb}{list}\PY{p}{(}\PY{n}{C\PYZus{}i}\PY{o}{.}\PY{n}{values}\PY{p}{(}\PY{p}{)}\PY{p}{)}\PY{p}{,}\PY{n}{bins}\PY{o}{=}\PY{n}{clu\PYZus{}bins}\PY{p}{,}\PY{n}{density}\PY{o}{=}\PY{k+kc}{True}\PY{p}{,}\PY{n}{color} \PY{o}{=} \PY{l+s+s1}{\PYZsq{}}\PY{l+s+s1}{b}\PY{l+s+s1}{\PYZsq{}}\PY{p}{,}\PY{n}{histtype}\PY{o}{=}\PY{l+s+s1}{\PYZsq{}}\PY{l+s+s1}{step}\PY{l+s+s1}{\PYZsq{}}\PY{p}{)}
\PY{n}{ax}\PY{p}{[}\PY{l+m+mi}{2}\PY{p}{]}\PY{o}{.}\PY{n}{loglog}\PY{p}{(}\PY{p}{)}
\PY{n}{ax}\PY{p}{[}\PY{l+m+mi}{2}\PY{p}{]}\PY{o}{.}\PY{n}{set\PYZus{}title}\PY{p}{(}\PY{l+s+s2}{\PYZdq{}}\PY{l+s+s2}{Clustering Distribution}\PY{l+s+s2}{\PYZdq{}}\PY{p}{)}
\PY{n}{ax}\PY{p}{[}\PY{l+m+mi}{2}\PY{p}{]}\PY{o}{.}\PY{n}{set\PYZus{}xlabel}\PY{p}{(}\PY{l+s+s2}{\PYZdq{}}\PY{l+s+s2}{Clustering Coefficient}\PY{l+s+s2}{\PYZdq{}}\PY{p}{)}
\PY{n}{ax}\PY{p}{[}\PY{l+m+mi}{2}\PY{p}{]}\PY{o}{.}\PY{n}{set\PYZus{}ylabel}\PY{p}{(}\PY{l+s+s2}{\PYZdq{}}\PY{l+s+s2}{Probability}\PY{l+s+s2}{\PYZdq{}}\PY{p}{)}
\PY{n}{ax}\PY{p}{[}\PY{l+m+mi}{2}\PY{p}{]}\PY{o}{.}\PY{n}{xaxis}\PY{o}{.}\PY{n}{set\PYZus{}inverted}\PY{p}{(}\PY{k+kc}{True}\PY{p}{)}
\PY{n}{ax}\PY{p}{[}\PY{l+m+mi}{2}\PY{p}{]}\PY{o}{.}\PY{n}{grid}\PY{p}{(}\PY{k+kc}{True}\PY{p}{)}

\PY{n}{fig}\PY{o}{.}\PY{n}{tight\PYZus{}layout}\PY{p}{(}\PY{p}{)}
\PY{n}{plt}\PY{o}{.}\PY{n}{show}\PY{p}{(}\PY{p}{)}
\PY{c+c1}{\PYZsh{}\PYZhy{}\PYZhy{}\PYZhy{}Configuration Model Version}
\PY{n}{my\PYZus{}null\PYZus{}model} \PY{o}{=} \PY{n}{configuration\PYZus{}model\PYZus{}from\PYZus{}degree\PYZus{}sequence}\PY{p}{(}\PY{n+nb}{list}\PY{p}{(}\PY{n}{deg}\PY{o}{.}\PY{n}{values}\PY{p}{(}\PY{p}{)}\PY{p}{)}\PY{p}{,}\PY{n}{return\PYZus{}simple}\PY{o}{=}\PY{k+kc}{True}\PY{p}{)}
\PY{n}{deg2} \PY{o}{=} \PY{n+nb}{dict}\PY{p}{(}\PY{n}{my\PYZus{}null\PYZus{}model}\PY{o}{.}\PY{n}{degree}\PY{p}{(}\PY{p}{)}\PY{p}{)}
\PY{n}{closeness\PYZus{}centrality2} \PY{o}{=} \PY{n}{nx}\PY{o}{.}\PY{n}{closeness\PYZus{}centrality}\PY{p}{(}\PY{n}{my\PYZus{}null\PYZus{}model}\PY{p}{)}
\PY{n}{C\PYZus{}i2} \PY{o}{=} \PY{n}{nx}\PY{o}{.}\PY{n}{clustering}\PY{p}{(}\PY{n}{my\PYZus{}null\PYZus{}model}\PY{p}{)}

\PY{n}{fig2}\PY{p}{,} \PY{n}{ax2} \PY{o}{=} \PY{n}{plt}\PY{o}{.}\PY{n}{subplots}\PY{p}{(}\PY{l+m+mi}{1}\PY{p}{,}\PY{l+m+mi}{3}\PY{p}{,}\PY{n}{figsize}\PY{o}{=}\PY{p}{(}\PY{l+m+mi}{8}\PY{p}{,}\PY{l+m+mf}{3.5}\PY{p}{)}\PY{p}{,}\PY{n}{dpi}\PY{o}{=}\PY{l+m+mi}{200}\PY{p}{)}
\PY{n}{deg\PYZus{}bins2} \PY{o}{=} \PY{n}{np}\PY{o}{.}\PY{n}{geomspace}\PY{p}{(}\PY{n}{np}\PY{o}{.}\PY{n}{min}\PY{p}{(}\PY{n+nb}{list}\PY{p}{(}\PY{n}{deg2}\PY{o}{.}\PY{n}{values}\PY{p}{(}\PY{p}{)}\PY{p}{)}\PY{p}{)}\PY{p}{,}\PY{n}{np}\PY{o}{.}\PY{n}{max}\PY{p}{(}\PY{n+nb}{list}\PY{p}{(}\PY{n}{deg2}\PY{o}{.}\PY{n}{values}\PY{p}{(}\PY{p}{)}\PY{p}{)}\PY{p}{)}\PY{p}{)}
\PY{n}{ax2}\PY{p}{[}\PY{l+m+mi}{0}\PY{p}{]}\PY{o}{.}\PY{n}{hist}\PY{p}{(}\PY{n+nb}{list}\PY{p}{(}\PY{n}{deg2}\PY{o}{.}\PY{n}{values}\PY{p}{(}\PY{p}{)}\PY{p}{)}\PY{p}{,}\PY{n}{bins}\PY{o}{=}\PY{n}{deg\PYZus{}bins2}\PY{p}{,} \PY{n}{density}\PY{o}{=}\PY{k+kc}{True}\PY{p}{,}\PY{n}{color} \PY{o}{=} \PY{l+s+s1}{\PYZsq{}}\PY{l+s+s1}{b}\PY{l+s+s1}{\PYZsq{}}\PY{p}{,}\PY{n}{histtype}\PY{o}{=}\PY{l+s+s1}{\PYZsq{}}\PY{l+s+s1}{step}\PY{l+s+s1}{\PYZsq{}}\PY{p}{)}
\PY{n}{ax2}\PY{p}{[}\PY{l+m+mi}{0}\PY{p}{]}\PY{o}{.}\PY{n}{loglog}\PY{p}{(}\PY{p}{)}
\PY{n}{ax2}\PY{p}{[}\PY{l+m+mi}{0}\PY{p}{]}\PY{o}{.}\PY{n}{set\PYZus{}title}\PY{p}{(}\PY{l+s+s2}{\PYZdq{}}\PY{l+s+s2}{Degree Distribution}\PY{l+s+s2}{\PYZdq{}}\PY{p}{)}
\PY{n}{ax2}\PY{p}{[}\PY{l+m+mi}{0}\PY{p}{]}\PY{o}{.}\PY{n}{set\PYZus{}xlabel}\PY{p}{(}\PY{l+s+s2}{\PYZdq{}}\PY{l+s+s2}{Degree}\PY{l+s+s2}{\PYZdq{}}\PY{p}{)}
\PY{n}{ax2}\PY{p}{[}\PY{l+m+mi}{0}\PY{p}{]}\PY{o}{.}\PY{n}{set\PYZus{}ylabel}\PY{p}{(}\PY{l+s+s2}{\PYZdq{}}\PY{l+s+s2}{Probability}\PY{l+s+s2}{\PYZdq{}}\PY{p}{)}
\PY{n}{ax2}\PY{p}{[}\PY{l+m+mi}{0}\PY{p}{]}\PY{o}{.}\PY{n}{grid}\PY{p}{(}\PY{k+kc}{True}\PY{p}{)}

\PY{n}{cen\PYZus{}bins2} \PY{o}{=} \PY{n}{np}\PY{o}{.}\PY{n}{geomspace}\PY{p}{(}\PY{n}{np}\PY{o}{.}\PY{n}{min}\PY{p}{(}\PY{n+nb}{list}\PY{p}{(}\PY{n}{closeness\PYZus{}centrality2}\PY{o}{.}\PY{n}{values}\PY{p}{(}\PY{p}{)}\PY{p}{)}\PY{p}{)}\PY{p}{,}\PY{n}{np}\PY{o}{.}\PY{n}{max}\PY{p}{(}\PY{n+nb}{list}\PY{p}{(}\PY{n}{closeness\PYZus{}centrality2}\PY{o}{.}\PY{n}{values}\PY{p}{(}\PY{p}{)}\PY{p}{)}\PY{p}{)}\PY{p}{)}
\PY{n}{ax2}\PY{p}{[}\PY{l+m+mi}{1}\PY{p}{]}\PY{o}{.}\PY{n}{hist}\PY{p}{(}\PY{n+nb}{list}\PY{p}{(}\PY{n}{closeness\PYZus{}centrality2}\PY{o}{.}\PY{n}{values}\PY{p}{(}\PY{p}{)}\PY{p}{)}\PY{p}{,}\PY{n}{bins}\PY{o}{=}\PY{n}{cen\PYZus{}bins2}\PY{p}{,}\PY{n}{density}\PY{o}{=}\PY{k+kc}{True}\PY{p}{,}\PY{n}{color} \PY{o}{=} \PY{l+s+s1}{\PYZsq{}}\PY{l+s+s1}{b}\PY{l+s+s1}{\PYZsq{}}\PY{p}{,}\PY{n}{histtype}\PY{o}{=}\PY{l+s+s1}{\PYZsq{}}\PY{l+s+s1}{step}\PY{l+s+s1}{\PYZsq{}}\PY{p}{)}
\PY{n}{ax2}\PY{p}{[}\PY{l+m+mi}{1}\PY{p}{]}\PY{o}{.}\PY{n}{loglog}\PY{p}{(}\PY{p}{)}
\PY{n}{ax2}\PY{p}{[}\PY{l+m+mi}{1}\PY{p}{]}\PY{o}{.}\PY{n}{set\PYZus{}title}\PY{p}{(}\PY{l+s+s2}{\PYZdq{}}\PY{l+s+s2}{Closeness Centrality Distribution}\PY{l+s+s2}{\PYZdq{}}\PY{p}{)}
\PY{n}{ax2}\PY{p}{[}\PY{l+m+mi}{1}\PY{p}{]}\PY{o}{.}\PY{n}{set\PYZus{}xlabel}\PY{p}{(}\PY{l+s+s2}{\PYZdq{}}\PY{l+s+s2}{Closeness Centrality}\PY{l+s+s2}{\PYZdq{}}\PY{p}{)}
\PY{n}{ax2}\PY{p}{[}\PY{l+m+mi}{1}\PY{p}{]}\PY{o}{.}\PY{n}{set\PYZus{}ylabel}\PY{p}{(}\PY{l+s+s2}{\PYZdq{}}\PY{l+s+s2}{Probability}\PY{l+s+s2}{\PYZdq{}}\PY{p}{)}
\PY{n}{ax2}\PY{p}{[}\PY{l+m+mi}{1}\PY{p}{]}\PY{o}{.}\PY{n}{xaxis}\PY{o}{.}\PY{n}{set\PYZus{}inverted}\PY{p}{(}\PY{k+kc}{True}\PY{p}{)}
\PY{n}{ax2}\PY{p}{[}\PY{l+m+mi}{1}\PY{p}{]}\PY{o}{.}\PY{n}{grid}\PY{p}{(}\PY{k+kc}{True}\PY{p}{)}

\PY{n}{clu\PYZus{}bins2} \PY{o}{=} \PY{n}{np}\PY{o}{.}\PY{n}{geomspace}\PY{p}{(}\PY{n}{np}\PY{o}{.}\PY{n}{min}\PY{p}{(}\PY{n+nb}{list}\PY{p}{(}\PY{n}{C\PYZus{}i2}\PY{o}{.}\PY{n}{values}\PY{p}{(}\PY{p}{)}\PY{p}{)}\PY{p}{)}\PY{o}{+}\PY{l+m+mf}{1e\PYZhy{}4}\PY{p}{,}\PY{n}{np}\PY{o}{.}\PY{n}{max}\PY{p}{(}\PY{n+nb}{list}\PY{p}{(}\PY{n}{C\PYZus{}i2}\PY{o}{.}\PY{n}{values}\PY{p}{(}\PY{p}{)}\PY{p}{)}\PY{p}{)}\PY{p}{)}
\PY{n}{ax2}\PY{p}{[}\PY{l+m+mi}{2}\PY{p}{]}\PY{o}{.}\PY{n}{hist}\PY{p}{(}\PY{n+nb}{list}\PY{p}{(}\PY{n}{C\PYZus{}i2}\PY{o}{.}\PY{n}{values}\PY{p}{(}\PY{p}{)}\PY{p}{)}\PY{p}{,}\PY{n}{bins}\PY{o}{=}\PY{n}{clu\PYZus{}bins2}\PY{p}{,}\PY{n}{density}\PY{o}{=}\PY{k+kc}{True}\PY{p}{,}\PY{n}{color} \PY{o}{=} \PY{l+s+s1}{\PYZsq{}}\PY{l+s+s1}{b}\PY{l+s+s1}{\PYZsq{}}\PY{p}{,}\PY{n}{histtype}\PY{o}{=}\PY{l+s+s1}{\PYZsq{}}\PY{l+s+s1}{step}\PY{l+s+s1}{\PYZsq{}}\PY{p}{)}
\PY{n}{ax2}\PY{p}{[}\PY{l+m+mi}{2}\PY{p}{]}\PY{o}{.}\PY{n}{loglog}\PY{p}{(}\PY{p}{)}
\PY{n}{ax2}\PY{p}{[}\PY{l+m+mi}{2}\PY{p}{]}\PY{o}{.}\PY{n}{set\PYZus{}title}\PY{p}{(}\PY{l+s+s2}{\PYZdq{}}\PY{l+s+s2}{Clustering Distribution}\PY{l+s+s2}{\PYZdq{}}\PY{p}{)}
\PY{n}{ax2}\PY{p}{[}\PY{l+m+mi}{2}\PY{p}{]}\PY{o}{.}\PY{n}{set\PYZus{}xlabel}\PY{p}{(}\PY{l+s+s2}{\PYZdq{}}\PY{l+s+s2}{Clustering Coefficient}\PY{l+s+s2}{\PYZdq{}}\PY{p}{)}
\PY{n}{ax2}\PY{p}{[}\PY{l+m+mi}{2}\PY{p}{]}\PY{o}{.}\PY{n}{set\PYZus{}ylabel}\PY{p}{(}\PY{l+s+s2}{\PYZdq{}}\PY{l+s+s2}{Probability}\PY{l+s+s2}{\PYZdq{}}\PY{p}{)}
\PY{n}{ax2}\PY{p}{[}\PY{l+m+mi}{2}\PY{p}{]}\PY{o}{.}\PY{n}{xaxis}\PY{o}{.}\PY{n}{set\PYZus{}inverted}\PY{p}{(}\PY{k+kc}{True}\PY{p}{)}
\PY{n}{ax2}\PY{p}{[}\PY{l+m+mi}{2}\PY{p}{]}\PY{o}{.}\PY{n}{grid}\PY{p}{(}\PY{k+kc}{True}\PY{p}{)}
\PY{n}{fig2}\PY{o}{.}\PY{n}{tight\PYZus{}layout}\PY{p}{(}\PY{p}{)}
\PY{n}{plt}\PY{o}{.}\PY{n}{show}\PY{p}{(}\PY{p}{)}
\end{Verbatim}
\end{tcolorbox}

    \begin{center}
    \adjustimage{max size={0.9\linewidth}{0.9\paperheight}}{Hughes_assignment01_files/Hughes_assignment01_41_0.png}
    \end{center}
    { \hspace*{\fill} \\}
    
    \begin{center}
    \adjustimage{max size={0.9\linewidth}{0.9\paperheight}}{Hughes_assignment01_files/Hughes_assignment01_41_1.png}
    \end{center}
    { \hspace*{\fill} \\}
    
    In my original dataset of friendships we can see the degree distribution
appears to follow a truncated power-law, with low-degree nodes occurring
commonly and high degree nodes occurring less and less commonly as
degree increases. Meanwhile, my closeness centrality appears to fall
into 2 buckets: very close or not very close, with both being on the
same order of magnitude in terms of probability to occur. Lastly, my
clustering coefficient is commonly high across all nodes, with a steep
drop off in probability of occurrence as soon as the clustering
coefficient begins to dip below 0.01.

In my configuration null model, we can see all these trends maintain
present albeit with some marginal differences. Namely, some of the
degree distribution between 10 and 100 appears slightly less probable in
the configuration model, and similarly high closeness centrality is a
little less probable than low closeness centrality (\emph{relative to
the original dataset}). Lastly, the clustering coefficient becomes most
common at a slightly lower value than the original dataset. This may
indicate that the original dataset was perhaps tighter knit relative to
the configuration model version, and thereby experienced slightly higher
degree, closeness and clustering as a result.

    \begin{enumerate}
\def\labelenumi{(\alph{enumi})}
\setcounter{enumi}{2}
\tightlist
\item
  Discuss the similarities and/or differences between your network
  dataset and the randomized versions of this network. How does this
  exercise---comparing network data to a null model---help us gain
  insights from network analysis?
\end{enumerate}

    The general trends in the datasets remained consistent: the degree
distribution resembling a truncated power law, the closeness centrality
falling into high and low buckets, and the clustering being generally
high. The main differences arose in the extent to which these trends
appeared, with the configuration model (i.e., randomized) version,
experiencing slightly lower probabilities for high centrality, high
clustering and high degreee (\emph{relatively speaking}).

Overall this exercise helps us gain insights into from networks
analysis, as it allows us to characterize how this network may be
different or unique from what is expected at random (\emph{even if it's
only a marginal that difference may be present}). It also allows us to
speculate and hypothesize how a network may be different, when we notice
a theme of marginal differences across different measures.

    \paragraph{(Optional) Question 5 (5 bonus
points).}\label{optional-question-5-5-bonus-points.}

For every problem set this semester, there will be an opportunity to
earn extra credit. Since this is a fairly new course (of a new masters
program), we are eager to receive genuine feedback on the material,
pace, instruction style, etc. of the course so far. Note: Learning how
to give effective feedback is an important skill during your PhD, so we
hope this is an opportunity for you to skill-build and improve the
course for future students. (Also, it's a bit odd to be evaluated on on
how good your feedback is, so we will be generous with awarding points
for these. With that said, please take some time to provide us with
productive comments!) If you need resources for giving effective
feedback, here are a few: *
https://hbr.org/2017/10/how-to-give-feedback-people-can-actually-use *
https://www.sessionlab.com/blog/feedback-techniques/

    I'm honestly really enjoying the course so far! I'm learning a lot and I
especially appreciate the in-class coding exercises we have been doing
lately. Despite working in data science, I've never been the strongest
coder, so I tend to take a fair amount of time to think and get through
problems, but it's because of that, that I think I feel that much more
confident in the material in the practice of it all. I think in that
sense the pro-active nature of the course has been a strong asset, in
that it was good there was the foresight to see that adding in exercises
could be worthwhile! I think similarly the capacity there has been to
allow lectures to speed-up or slow-down, ask questions, and go on
tangents and really explore and play with the material has added a lot
of value and is very appreciated! I can see a lot of students joining
and taking this course to study a specific kind of network, or coming in
with a different background and wanting or needing to know more about
some topics over others.

I'm afraid I don't have too much constructive criticism for the course
at the moment. I think if I had to name some it would be, that I've
noticed that the students outside of my NetSci courses (\emph{right now
I'm taking 1 other class called Big Data for Cities that has a lot of
MPP and Structural Engineering kinds of students}) tended to have little
to no background in coding, so I could see in the future once we get
more students, we may see students with less background struggling. That
all being said, I also realize that at the moment we are a small cohort
and as such there isn't really a huge needed to go slow like that. In
any case, I mention it just as something to be cognizant of in case it's
something the course encounters in the future.

Other than that, and this is perhaps nitpick territory, but it could be
valuable to have any windows specific coding errors mentioned perhaps as
like a little `heads-up' in some of the code, but it honestly hasn't
really been a problem at all in my experience. I think I had 1 error in
reading a file due to Windows but I was able to just tell Python to
ignore it as an argument to the read() and then it worked fine.

Hope this feedback is of help, and looking forward to the rest of the
course!


    % Add a bibliography block to the postdoc
    
    
    
\end{document}
